%%
%% @Author           : Albert Wang
%% @Time             : 2023-01-08 21:21:02
%% @Description      :
%% @Email            : shadowofgost@outlook.com
%% @FilePath         : /BachelorThesis/main.tex
%% @LastTime         : 2024-03-29 09:29:05
%% @LastAuthor       : Albert Wang
%% @Software         : Vscode
%% @ Copyright Notice : Copyright (c) 2023 Albert Wang 王子睿, All Rights Reserved.
%%
% !TeX encoding = UTF-8
% !TeX program = xelatex
% !TeX spellcheck = en_US
% !TEX root = /workspaces/BachelorThesis/main.tex.
\documentclass[degree=bachelor,print-style=oneside,sequence-style=liberal]{sudathesis}
  % 学位 degree:
  %   doctor | master | bachelor | postdoc
  % 学位类型 degree-type:
  %   academic（默认）| professional
  % 语言 language
  %   chinese（默认）| english
  % 字体库 fontset
  %   windows | mac | fandol | ubuntu
  % 建议终版使用 Windows 平台的字体编译
  % 打印设置 print-style
  %    oneside | twoside
  % 章节编码设置 sequence-style
  %   science | liberal 默认science
  % 默认使用双面打印，根据自己需要设定双面还是单面打印，如果是提交电子版论文建议使用单面打印消除空白页


%%% 配置setup.tex 和main。tex后运行，记得自己选择导入生成封面和选择参考文献格式


% 论文基本配置，加载宏包等全局配置
%%
%% @Author           : Albert Wang
%% @Time             : 2023-01-08 21:21:02
%% @Description      :
%% @Email            : shadowofgost@outlook.com
%% @FilePath         : /1-AccountingThesis/BachelorThesis/setup.tex
%% @LastTime         : 2024-03-31 22:27:05
%% @LastAuthor       : Albert Wang
%% @Software         : Vscode
%% @ Copyright Notice : Copyright (c) 2023 Albert Wang 王子睿, All Rights Reserved.
%%

% !TeX root = ./main.tex

% 论文基本信息配置

\thusetup{
  %******************************
  % 注意：
  %   1. 配置里面不要出现空行
  %   2. 不需要的配置信息可以删除
  %   3. 建议先阅读文档中所有关于选项的说明
  %******************************
  %
  % 输出格式
  %   选择打印版（print）或用于提交的电子版（electronic），前者会插入空白页以便直接双面打印
  %
  output = print,
  %
  % 标题
  %   可使用“\\”命令手动控制换行
  %
  title  = {我国新能源汽车上市公司最优资本结构研究——以比亚迪公司为例},
  title* = {Research on the optimal capital structure of listed new energy vehicle companies in China: A case study of BYD Company.},
  % 副标题
  subtitle  = {这里是副标题},
  subtitle* = {Here is Subtitle},
  %
  % 学科门类
  %   1. 学术型
  %      - 中文
  %        需注明所属的学科门类，例如：
  %        哲学、经济学、法学、教育学、文学、历史学、理学、工学、农学、医学、
  %        军事学、管理学、艺术学
  %      - 英文
  %        博士：Doctor of Philosophy
  %        硕士：
  %          哲学、文学、历史学、法学、教育学、艺术学门类，公共管理学科
  %          填写“Master of Arts“，其它填写“Master of Science”
  %   2. 专业型
  %      直接填写专业学位的名称，例如：
  %      教育博士、工程硕士等
  %      Doctor of Education, Master of Engineering
  %   3. 本科生不需要填写
  %
  degree-category  = {工学硕士},
  degree-category* = {Master of Science},
  %
  % 培养单位
  %   填写所属院系的全名
  %
  department = {会计系},
  department* = {Accountant Departmen},
  %
  % 学科
  %   1. 研究生学术型学位，获得一级学科授权的学科填写一级学科名称，其他填写二级学科名称
  %   2. 本科生填写专业名称，第二学位论文需标注“（第二学位）”
  %
  discipline  = {计算机科学与技术},
  discipline* = {Computer Science and Technology},
  %
  % 学科研究方向%
  %
  discipline-direct = {软件工程},
  discipline-direct* = {Computer Science and Technology},
  %
  % 专业领域
  %   1. 设置专业领域的专业学位类别，填写相应专业领域名称
  %   2. 2019 级及之前工程硕士学位论文，在 `engineering-field` 填写相应工程领域名称
  %   3. 其他专业学位类别的学位论文无需此信息
  %
  % professional-field  = {计算机技术},
  % professional-field* = {Computer Technology},
  %
  % 姓名
  %
  author  = {王子睿},
  author* = {Wang Zirui},
  % 学号
  student-id={1910404051},
  %
  % 指导教师
  %   中文姓名和职称之间以英文逗号“,”分开，下同
  %
  supervisor  = {龚蕾, 副教授},
  supervisor* = {Associate Professor Gong Lei},
  %
  % 副指导教师
  %
  associate-supervisor  = {陈文光, 教授},
  associate-supervisor* = {Professor Chen Wenguang},
  %
  % 联合指导教师
  %
  % co-supervisor  = {某某某, 教授},
  % co-supervisor* = {Professor Mou Moumou},
  %
  % 日期
  %   使用 ISO 格式；默认为当前时间
  %
  % date = {2019-07-07},
  %
  % 是否在中文封面后的空白页生成书脊（默认 false）
  %
  include-spine = false,
  %
  % 密级和年限
  %   秘密, 机密, 绝密
  %
  % secret-level = {秘密},
  % secret-year  = {10},
  %
  % 博士后专有部分
  %
  % clc                = {分类号},
  % udc                = {UDC},
  % id                 = {编号},
  % discipline-level-1 = {计算机科学与技术},  % 流动站（一级学科）名称
  % discipline-level-2 = {系统结构},          % 专业（二级学科）名称
  % start-date         = {2011-07-01},        % 研究工作起始时间
  % end-date           = {2011-07-01},        % 研究工作结束时间
}

% 载入所需的宏包

% 定理类环境宏包
\usepackage{amsthm}
% 也可以使用 ntheorem
% \usepackage[amsmath,thmmarks,hyperref]{ntheorem}

\thusetup{
  %
  % 数学字体
  % math-style = GB,  % GB | ISO | TeX
  math-font  = xits,  % stix | xits | libertinus
}

% 可以使用 nomencl 生成符号和缩略语说明
% \usepackage{nomencl}
% \makenomenclature

% 表格加脚注
\usepackage{threeparttable}

% 表格中支持跨行
\usepackage{multirow}

% 固定宽度的表格。
% \usepackage{tabularx}

% 跨页表格
\usepackage{longtable}

% 算法
\usepackage{algorithm}
\usepackage{algorithmic}

% 量和单位
\usepackage{siunitx}

% 序列编号的方式
\usepackage{enumitem}
% 参考文献使用 BibTeX + natbib 宏包
% 顺序编码制
%\usepackage[sort]{natbib}
%\bibliographystyle{./bibstyle/thuthesis-numeric}

% 著者-出版年制
\usepackage[sort]{natbib}
\bibliographystyle{./bibstyle/thuthesis-author-year}
\thusetup{
  cite-style = super,
}
% 本科生参考文献的著录格式
%\usepackage[sort]{natbib}
%\bibliographystyle{./bibstyle/thuthesis-bachelor}

% 参考文献使用 BibLaTeX 宏包
% \usepackage[style=./bibstyle/thuthesis-numeric]{biblatex}
% \usepackage[style=./bibstyle/thuthesis-author-year]{biblatex}
% \usepackage[style=apa]{biblatex}
% \usepackage[style=mla-new]{biblatex}
% 声明 BibLaTeX 的数据库
% \addbibresource{reference/AccountThesis.bib}

% 定义所有的图片文件在 figures 子目录下
\graphicspath{{figures/}}

% 数学命令
\makeatletter
\newcommand\dif{%  % 微分符号
  \mathop{}\!%
  \ifthu@math@style@TeX
    d%
  \else
    \mathrm{d}%
  \fi
}
\makeatother

% hyperref 宏包在最后调用
\usepackage{hyperref}


\begin{document}

% 封面，使用清华原版的封面
%\maketitle
% 插入封面，针对苏大本科生毕业论文
\includepdf[pages={1}]{./figures/cover-zh.pdf}
% 生成封面，针对苏大研究生毕业论文中文版
%\titlech
% 生成封面，针对苏大研究生毕业论文英文版
%\titleeng
% 学位论文指导小组、公开评阅人和答辩委员会名单
% 本科生不需要
%%%
%% @Author           : Albert Wang
%% @Time             : 2023-01-08 21:21:02
%% @Description      :
%% @Email            : shadowofgost@outlook.com
%% @FilePath         : /Sudathesis-Bachelor-Template/contents/committee.tex
%% @LastTime         : 2023-01-10 11:11:25
%% @LastAuthor       : Albert Wang
%% @Software         : Vscode
%% @ Copyright Notice : Copyright (c) 2023 Albert Wang 王子睿, All Rights Reserved.
%%

% !TeX root = ../main.tex

\begin{committee}[name={学位论文指导小组、公开评阅人和答辩委员会名单}]

  \newcolumntype{C}[1]{@{}>{\centering\arraybackslash}p{#1}}

  \section*{指导小组名单}

  \begin{center}
    \begin{tabular}{C{3cm}C{3cm}C{9cm}@{}}
      李XX & 教授     & 清华大学 \\
      王XX & 副教授   & 清华大学 \\
      张XX & 助理教授 & 清华大学 \\
    \end{tabular}
  \end{center}


  \section*{公开评阅人名单}

  \begin{center}
    \begin{tabular}{C{3cm}C{3cm}C{9cm}@{}}
      刘XX & 教授   & 清华大学                    \\
      陈XX & 副教授 & XXXX大学                    \\
      杨XX & 研究员 & 中国XXXX科学院XXXXXXX研究所 \\
    \end{tabular}
  \end{center}


  \section*{答辩委员会名单}

  \begin{center}
    \begin{tabular}{C{2.75cm}C{2.98cm}C{4.63cm}C{4.63cm}@{}}
      主席 & 赵XX                  & 教授                    & 清华大学       \\
      委员 & 刘XX                  & 教授                    & 清华大学       \\
          & \multirow{2}{*}{杨XX} & \multirow{2}{*}{研究员} & 中国XXXX科学院 \\
          &                       &                         & XXXXXXX研究所  \\
          & 黄XX                  & 教授                    & XXXX大学       \\
          & 周XX                  & 副教授                  & XXXX大学       \\
      秘书 & 吴XX                  & 助理研究员              & 清华大学       \\
    \end{tabular}
  \end{center}

\end{committee}



% 也可以导入 Word 版转的 PDF 文件
% \begin{committee}[file=figures/committee.pdf]
% \end{committee}


% 使用授权的说明
%\copyrightpage
% 将签字扫描后授权文件 scan-copyright.pdf 替换原始页面
% \copyrightpage[file=scan-copyright.pdf]

% 插入独创性声明和使用授权说明（苏大本科生）
\includepdf[pages={1,2}]{./figures/statement.pdf}

\frontmatter
%\newpage
%\pagestyle{cnfrontmatter}
%%%
%% @Author           : Albert Wang
%% @Time             : 2023-01-08 21:21:02
%% @Description      :
%% @Email            : shadowofgost@outlook.com
%% @FilePath         : /Sudathesis-Bachelor-Template/contents/abstract.tex
%% @LastTime         : 2023-01-10 11:07:23
%% @LastAuthor       : Albert Wang
%% @Software         : Vscode
%% @ Copyright Notice : Copyright (c) 2023 Albert Wang 王子睿, All Rights Reserved.
%%

% !TeX root = ../main.tex

% 中英文摘要和关键字

\begin{abstract}
  论文的摘要是对论文研究内容和成果的高度概括。
  摘要应对论文所研究的问题及其研究目的进行描述，对研究方法和过程进行简单介绍，对研究成果和所得结论进行概括。
  摘要应具有独立性和自明性，其内容应包含与论文全文同等量的主要信息。
  使读者即使不阅读全文，通过摘要就能了解论文的总体内容和主要成果。

  论文摘要的书写应力求精确、简明。
  切忌写成对论文书写内容进行提要的形式，尤其要避免“第 1 章……；第 2 章……；……”这种或类似的陈述方式。

  关键词是为了文献标引工作、用以表示全文主要内容信息的单词或术语。
  关键词不超过 5 个，每个关键词中间用分号分隔。

  % 关键词用“英文逗号”分隔，输出时会自动处理为正确的分隔符
  \thusetup{
    keywords = {关键词 1, 关键词 2, 关键词 3, 关键词 4, 关键词 5},
  }
\end{abstract}

%\newpage
%\pagestyle{enfrontmatter}
%%%
%% @Author           : Albert Wang
%% @Time             : 2023-05-11 15:10:33
%% @Description      :
%% @Email            : shadowofgost@outlook.com
%% @FilePath         : /Sudathesis-Bachelor-Template/contents/abstract_en.tex
%% @LastTime         : 2023-05-11 15:11:07
%% @LastAuthor       : Albert Wang
%% @Software         : Vscode
%% @ Copyright Notice : Copyright (c) 2023 Albert Wang 王子睿, All Rights Reserved.
%%

\begin{abstract*}
  An abstract of a dissertation is a summary and extraction of research work and contributions.
  Included in an abstract should be description of research topic and research objective, brief introduction to methodology and research process, and summary of conclusion and contributions of the research.
  An abstract should be characterized by independence and clarity and carry identical information with the dissertation.
  It should be such that the general idea and major contributions of the dissertation are conveyed without reading the dissertation.

  An abstract should be concise and to the point.
  It is a misunderstanding to make an abstract an outline of the dissertation and words “the first chapter”, “the second chapter” and the like should be avoided in the abstract.

  Keywords are terms used in a dissertation for indexing, reflecting core information of the dissertation.
  An abstract may contain a maximum of 5 keywords, with semi-colons used in between to separate one another.

  % Use comma as separator when inputting
  \thusetup{
    keywords* = {keyword 1, keyword 2, keyword 3, keyword 4, keyword 5},
  }
\end{abstract*}

%\newpage
%\pagestyle{plain}
%\frontmatter
% !TeX root = ../thuthesis-example.tex

% 中英文摘要和关键字

\begin{abstract}
当前，随着新一轮科技和产业变革的到来，尽管新能源汽车得到了全球推崇，但近期新能源汽车补贴政策的逐步减少也带来了一定压力。在汽车行业，新能源汽车所需的先进技术和资金支持变得尤为重要。随着各类国产新能源厂商的崛起，新能源企业之间的竞争越来越激烈。并且逐渐呈现出价格战的趋势。因此，从长远来看，搭建模型优化新能源企业的资本结构使其达到最佳的状态对企业而言具有重要意义。

本文以比亚迪公司为研究对象，重点分析介绍了比亚迪目前的资本结构和财务情况，通过横向和纵向比较，全面剖析了比亚迪公司目前的资本结构，主要的资金占比，资金来源和融资情况对企业发展的支撑作用。

在这基础上通过从融资方式、债务期限和债务结构等角度入手，揭示了比亚迪公司存在的问题，包括偏重债务融资、债务期限短、直接融资规模受限、较少采用股权融资等。这些问题的核心原因在于未能科学评估融资需求，短期债务对业务推动更明显和受限的资本市场发展。

最后，本文基于静态方程和熵权法分别建立静态优化模型和动态优化模型，计算出比亚迪公司近几年的最优资本结构。并从加强营运资金管理、合理分配债务期限结构和合理调整企业融资渠道比重三个方面针对比亚迪资本结构现状提出优化的建议。 
%本文以比亚迪公司为研究对象，基于资本结构静态优化理论和动态优化理论还有对以往学者对企业最优资本结构的研究，对企业当前资本结构状态进行了深入的分析和研究，并基于静态优化公式和熵权法搭建了资本结构优化模型，为新能源汽车企业的资本结构优化提供了参考。

%首先，通过微观角度深入分析比亚迪公司的资本结构现状。本文分析了比亚迪公司的发展状况以及新能源汽车行业的前景，认为比亚迪公司在新能源汽车领域发展迅速且前景广阔。同时，重点分析介绍了比亚迪公司目前的资本结构和财务情况，通过横向和纵向比较，全面剖析了比亚迪公司目前的资本结构，主要的资金占比，资金来源和融资情况对企业发展的支撑作用。

%其次，本文深入分析了比亚迪公司现有的资本结构存在的问题并探讨了其原因。通过从融资方式、债务期限和债务结构等角度入手，揭示了比亚迪公司存在的问题，包括偏重债务融资、债务期限短、直接融资规模受限、较少采用股权融资等。这些问题的核心原因在于未能科学评估融资需求，短期债务对业务推动更明显和受限的资本市场发展。

%最后，本文提出了解决比亚迪资本结构问题的优化策略。结合比亚迪公司的实际情况，明确了优化融资结构的目标。根据预测结果，通过优化路径，探索了比亚迪公司的最佳资本结构，并提出了针对性的优化策略，以期为比亚迪公司的未来发展提供有效支持。
  \thusetup{
    keywords = {比亚迪公司, 新能源汽车, 资本结构, 优化策略, 熵权法},
  }
\end{abstract}

\begin{abstract*}
With the advent of a new wave of technological and industrial revolution, global recognition of electric vehicles (EVs) has surged. However, the gradual reduction in subsidies for EVs has also created pressure. In the automotive industry, advanced technology and financial support are crucial for the development of new energy vehicles. As domestic new energy manufacturers rise, competition among these companies intensifies, leading to a trend of price wars. Therefore, from a long-term perspective, optimizing the capital structure of new energy enterprises is essential to achieve their optimal state.

This study focuses on BYD Company as the research subject. We analyze and introduce BYD’s current capital structure and financial situation. Through both horizontal and vertical comparisons, we comprehensively dissect BYD’s capital structure, major funding sources, and financing conditions that support its growth.

Building upon this analysis, we delve into issues faced by BYD, including an overreliance on debt financing, short debt maturity, limited direct financing scale, and underutilization of equity financing. The core reason behind these challenges lies in the inadequate assessment of financing needs, with short-term debt significantly impacting business expansion and restricting capital market development.

Finally, we establish static and dynamic optimization models using static equations and entropy weight methods. These models calculate BYD’s optimal capital structure over recent years. Based on our findings, we propose recommendations to enhance working capital management, allocate debt maturity structures rationally, and adjust the proportion of financing channels for BYD’s current capital structure.
  \thusetup{
    keywords* = {BYD company, new energy vehicles, capital structure; optimization plan, entropy weight method},
  }
\end{abstract*}

% 目录
\tableofcontents

% 插图和附表清单
% 本科生的插图索引和表格索引需要移至正文之后、参考文献前
% \listoffiguresandtables  % 插图和附表清单（仅限研究生）
%\listoffigures           % 插图清单
%\listoftables            % 附表清单

% 符号对照表
% %%
%% @Author           : Albert Wang
%% @Time             : 2023-01-08 21:21:02
%% @Description      :
%% @Email            : shadowofgost@outlook.com
%% @FilePath         : /Sudathesis-Bachelor-Template/contents/denotation.tex
%% @LastTime         : 2023-01-10 11:11:31
%% @LastAuthor       : Albert Wang
%% @Software         : Vscode
%% @ Copyright Notice : Copyright (c) 2023 Albert Wang 王子睿, All Rights Reserved.
%%

% !TeX root = ../main.tex

\begin{denotation}[3cm]
  \item[PI] 聚酰亚胺
  \item[MPI] 聚酰亚胺模型化合物，N-苯基邻苯酰亚胺
  \item[PBI] 聚苯并咪唑
  \item[MPBI] 聚苯并咪唑模型化合物，N-苯基苯并咪唑
  \item[PY] 聚吡咙
  \item[PMDA-BDA] 均苯四酸二酐与联苯四胺合成的聚吡咙薄膜
  \item[MPY] 聚吡咙模型化合物
  \item[As-PPT] 聚苯基不对称三嗪
  \item[MAsPPT] 聚苯基不对称三嗪单模型化合物，3,5,6-三苯基-1,2,4-三嗪
  \item[DMAsPPT] 聚苯基不对称三嗪双模型化合物（水解实验模型化合物）
  \item[S-PPT] 聚苯基对称三嗪
  \item[MSPPT] 聚苯基对称三嗪模型化合物，2,4,6-三苯基-1,3,5-三嗪
  \item[PPQ] 聚苯基喹噁啉
  \item[MPPQ] 聚苯基喹噁啉模型化合物，3,4-二苯基苯并二嗪
  \item[HMPI] 聚酰亚胺模型化合物的质子化产物
  \item[HMPY] 聚吡咙模型化合物的质子化产物
  \item[HMPBI] 聚苯并咪唑模型化合物的质子化产物
  \item[HMAsPPT] 聚苯基不对称三嗪模型化合物的质子化产物
  \item[HMSPPT] 聚苯基对称三嗪模型化合物的质子化产物
  \item[HMPPQ] 聚苯基喹噁啉模型化合物的质子化产物
  \item[PDT] 热分解温度
  \item[HPLC] 高效液相色谱（High Performance Liquid Chromatography）
  \item[HPCE] 高效毛细管电泳色谱（High Performance Capillary lectrophoresis）
  \item[LC-MS] 液相色谱-质谱联用（Liquid chromatography-Mass Spectrum）
  \item[TIC] 总离子浓度（Total Ion Content）
  \item[\textit{ab initio}] 基于第一原理的量子化学计算方法，常称从头算法
  \item[DFT] 密度泛函理论（Density Functional Theory）
  \item[$E_a$] 化学反应的活化能（Activation Energy）
  \item[ZPE] 零点振动能（Zero Vibration Energy）
  \item[PES] 势能面（Potential Energy Surface）
  \item[TS] 过渡态（Transition State）
  \item[TST] 过渡态理论（Transition State Theory）
  \item[$\increment G^\neq$] 活化自由能（Activation Free Energy）
  \item[$\kappa$] 传输系数（Transmission Coefficient）
  \item[IRC] 内禀反应坐标（Intrinsic Reaction Coordinates）
  \item[$\nu_i$] 虚频（Imaginary Frequency）
  \item[ONIOM] 分层算法（Our own N-layered Integrated molecular Orbital and molecular Mechanics）
  \item[SCF] 自洽场（Self-Consistent Field）
  \item[SCRF] 自洽反应场（Self-Consistent Reaction Field）
\end{denotation}



% 也可以使用 nomencl 宏包，需要在导言区
% \usepackage{nomencl}
% \makenomenclature

% 在这里输出符号说明
% \printnomenclature[3cm]

% 在正文中的任意为都可以标题
% \nomenclature{PI}{聚酰亚胺}
% \nomenclature{MPI}{聚酰亚胺模型化合物，N-苯基邻苯酰亚胺}
% \nomenclature{PBI}{聚苯并咪唑}
% \nomenclature{MPBI}{聚苯并咪唑模型化合物，N-苯基苯并咪唑}
% \nomenclature{PY}{聚吡咙}
% \nomenclature{PMDA-BDA}{均苯四酸二酐与联苯四胺合成的聚吡咙薄膜}
% \nomenclature{MPY}{聚吡咙模型化合物}
% \nomenclature{As-PPT}{聚苯基不对称三嗪}
% \nomenclature{MAsPPT}{聚苯基不对称三嗪单模型化合物，3,5,6-三苯基-1,2,4-三嗪}
% \nomenclature{DMAsPPT}{聚苯基不对称三嗪双模型化合物（水解实验模型化合物）}
% \nomenclature{S-PPT}{聚苯基对称三嗪}
% \nomenclature{MSPPT}{聚苯基对称三嗪模型化合物，2,4,6-三苯基-1,3,5-三嗪}
% \nomenclature{PPQ}{聚苯基喹噁啉}
% \nomenclature{MPPQ}{聚苯基喹噁啉模型化合物，3,4-二苯基苯并二嗪}
% \nomenclature{HMPI}{聚酰亚胺模型化合物的质子化产物}
% \nomenclature{HMPY}{聚吡咙模型化合物的质子化产物}
% \nomenclature{HMPBI}{聚苯并咪唑模型化合物的质子化产物}
% \nomenclature{HMAsPPT}{聚苯基不对称三嗪模型化合物的质子化产物}
% \nomenclature{HMSPPT}{聚苯基对称三嗪模型化合物的质子化产物}
% \nomenclature{HMPPQ}{聚苯基喹噁啉模型化合物的质子化产物}
% \nomenclature{PDT}{热分解温度}
% \nomenclature{HPLC}{高效液相色谱（High Performance Liquid Chromatography）}
% \nomenclature{HPCE}{高效毛细管电泳色谱（High Performance Capillary lectrophoresis）}
% \nomenclature{LC-MS}{液相色谱-质谱联用（Liquid chromatography-Mass Spectrum）}
% \nomenclature{TIC}{总离子浓度（Total Ion Content）}
% \nomenclature{\textit{ab initio}}{基于第一原理的量子化学计算方法，常称从头算法}
% \nomenclature{DFT}{密度泛函理论（Density Functional Theory）}
% \nomenclature{$E_a$}{化学反应的活化能（Activation Energy）}
% \nomenclature{ZPE}{零点振动能（Zero Vibration Energy）}
% \nomenclature{PES}{势能面（Potential Energy Surface）}
% \nomenclature{TS}{过渡态（Transition State）}
% \nomenclature{TST}{过渡态理论（Transition State Theory）}
% \nomenclature{$\increment G^\neq$}{活化自由能（Activation Free Energy）}
% \nomenclature{$\kappa$}{传输系数（Transmission Coefficient）}
% \nomenclature{IRC}{内禀反应坐标（Intrinsic Reaction Coordinates）}
% \nomenclature{$\nu_i$}{虚频（Imaginary Frequency）}
% \nomenclature{ONIOM}{分层算法（Our own N-layered Integrated molecular Orbital and molecular Mechanics）}
% \nomenclature{SCF}{自洽场（Self-Consistent Field）}
% \nomenclature{SCRF}{自洽反应场（Self-Consistent Reaction Field）}


%\newpage
% 正文部分
\mainmatter
%%%
%% @Author           : Albert Wang
%% @Time             : 2023-01-08 21:21:02
%% @Description      :
%% @Email            : shadowofgost@outlook.com
%% @FilePath         : /1-AccountingThesis/BachelorThesis/contents/chap01.tex
%% @LastTime         : 2023-12-25 21:23:19
%% @LastAuthor       : Albert Wang
%% @Software         : Vscode
%% @ Copyright Notice : Copyright (c) 2023 Albert Wang 王子睿, All Rights Reserved.
%%

% !TeX root = ../main.tex

\chapter{论文主要部分的写法}

研究生学位论文撰写，除表达形式上需要符合一定的格式要求外，内容方面上也要遵循一些共性原则。

通常研究生学位论文只能有一个主题（不能是几块工作拼凑在一起），该主题应针对某学科领域中的一个具体问题展开深入、系统的研究，并得出有价值的研究结论。
学位论文的研究主题切忌过大，例如，“中国国有企业改制问题研究”这样的研究主题过大，因为“国企改制”涉及的问题范围太广，很难在一本研究生学位论文中完全研究透彻。



\section{论文的语言及表述}

除国际研究生外，学位论文一律须用汉语书写。
学位论文应当用规范汉字进行撰写，除古汉语研究中涉及的古文字和参考文献中引用的外文文献之外，均采用简体汉字撰写。

国际研究生一般应以中文或英文书写学位论文，格式要求同上。
论文须用中文封面。

研究生学位论文是学术作品，因此其表述要严谨简明，重点突出，专业常识应简写或不写，做到立论正确、数据可靠、说明透彻、推理严谨、文字凝练、层次分明，避免使用文学性质的或带感情色彩的非学术性语言。

论文中如出现一个非通用性的新名词、新术语或新概念，需随即解释清楚。



\section{论文题目的写法}

论文题目应简明扼要地反映论文工作的主要内容，力求精炼、准确，切忌笼统。
论文题目是对研究对象的准确、具体描述，一般要在一定程度上体现研究结论，因此，论文题目不仅应告诉读者这本论文研究了什么问题，更要告诉读者这个研究得出的结论。
例如：“在事实与虚构之间：梅乐、卡彭特、沃尔夫的新闻观”就比“三个美国作家的新闻观研究”更专业、更准确。



\section{摘要的写法}

论文摘要是对论文研究内容的高度概括，应具有独立性和自含性，即应是 一篇简短但意义完整的文章。
通过阅读论文摘要，读者应该能够对论文的研究 方法及结论有一个整体性的了解，因此摘要的写法应力求精确简明。
论文摘要 应包括对问题及研究目的的描述、对使用的方法和研究过程进行的简要介绍、 对研究结论的高度凝练等，重点是结果和结论。

论文摘要切忌写成全文的提纲，尤其要避免“第 1 章……；第 2 章……；……”这样的陈述方式。



\section{引言的写法}

一篇学位论文的引言大致包含如下几个部分：
1、问题的提出；
2、选题背 景及意义；
3、文献综述；
4、研究方法；
5、论文结构安排。
\begin{itemize}
  \item 问题的提出：要清晰地阐述所要研究的问题“是什么”。
    \footnote{选题时切记要有“问题意识”，不要选不是问题的问题来研究。}
  \item 选题背景及意义：论述清楚为什么选择这个题目来研究，即阐述该研究对学科发展的贡献、对国计民生的理论与现实意义等。
  \item 文献综述：对本研究主题范围内的文献进行详尽的综合述评，“述”的同时一定要有“评”，指出现有研究状态，仍存在哪些尚待解决的问题，讲出自己的研究有哪些探索性内容。
  \item 研究方法：讲清论文所使用的学术研究方法。
  \item 论文结构安排：介绍本论文的写作结构安排。
\end{itemize}



\section{正文的写法}

本部分是论文作者的研究内容，不能将他人研究成果不加区分地掺和进来。
已经在引言的文献综述部分讲过的内容，这里不需要再重复。
各章之间要存在有机联系，符合逻辑顺序。



\section{结论的写法}

结论是对论文主要研究结果、论点的提炼与概括，应精炼、准确、完整，使读者看后能全面了解论文的意义、目的和工作内容。
结论是最终的、总体的结论，不是正文各章小结的简单重复。
结论应包括论文的核心观点，主要阐述作者的创造性工作及所取得的研究成果在本领域中的地位、作用和意义，交代研究工作的局限，提出未来工作的意见或建议。
同时，要严格区分自己取得的成果与指导教师及他人的学术成果。

在评价自己的研究工作成果时，要实事求是，除非有足够的证据表明自己的研究是“首次”、“领先”、“填补空白”的，否则应避免使用这些或类似词语。


%%%
%% @Author           : Albert Wang
%% @Time             : 2023-01-08 21:21:02
%% @Description      :
%% @Email            : shadowofgost@outlook.com
%% @FilePath         : /Sudathesis-Bachelor-Template/contents/chap02.tex
%% @LastTime         : 2023-01-10 11:10:53
%% @LastAuthor       : Albert Wang
%% @Software         : Vscode
%% @ Copyright Notice : Copyright (c) 2023 Albert Wang 王子睿, All Rights Reserved.
%%

% !TeX root = ../main.tex

\chapter{图表示例}

\section{插图}

图片通常在 \env{figure} 环境中使用 \cs{includegraphics} 插入，如图~\ref{fig:example} 的源代码。
建议矢量图片使用 PDF 格式，比如数据可视化的绘图；
照片应使用 JPG 格式；
其他的栅格图应使用无损的 PNG 格式。
注意，LaTeX 不支持 TIFF 格式；EPS 格式已经过时。

\begin{figure}
  \centering
  \includegraphics[width=0.5\linewidth]{example-image-a.pdf}
  \caption*{国外的期刊习惯将图表的标题和说明文字写成一段，需要改写为标题只含图表的名称，其他说明文字以注释方式写在图表下方，或者写在正文中。}
  \caption{示例图片标题}
  \label{fig:example}
\end{figure}

若图或表中有附注，采用英文小写字母顺序编号，附注写在图或表的下方。
国外的期刊习惯将图表的标题和说明文字写成一段，需要改写为标题只含图表的名称，其他说明文字以注释方式写在图表下方，或者写在正文中。

如果一个图由两个或两个以上分图组成时，各分图分别以 (a)、(b)、(c)...... 作为图序，并须有分图题。
推荐使用 \pkg{subcaption} 宏包来处理， 比如图~\ref{fig:subfig-a} 和图~\ref{fig:subfig-b}。

\begin{figure}
  \centering
  \subcaptionbox{分图 A\label{fig:subfig-a}}
    {\includegraphics[width=0.35\linewidth]{example-image-a.pdf}}
  \subcaptionbox{分图 B\label{fig:subfig-b}}
    {\includegraphics[width=0.35\linewidth]{example-image-b.pdf}}
  \caption{多个分图的示例}
  \label{fig:multi-image}
\end{figure}



\section{表格}

表应具有自明性。为使表格简洁易读，尽可能采用三线表，如表~\ref{tab:three-line}。
三条线可以使用 \pkg{booktabs} 宏包提供的命令生成。

\begin{table}
  \centering
  \caption{三线表示例}
  \begin{tabular}{ll}
    \toprule
    文件名          & 描述                         \\
    \midrule
    thuthesis.dtx   & 模板的源文件，包括文档和注释 \\
    thuthesis.cls   & 模板文件                     \\
    thuthesis-*.bst & BibTeX 参考文献表样式文件    \\
    \bottomrule
  \end{tabular}
  \label{tab:three-line}
\end{table}

表格如果有附注，尤其是需要在表格中进行标注时，可以使用 \pkg{threeparttable} 宏包。
研究生要求使用英文小写字母 a、b、c……顺序编号，本科生使用圈码 ①、②、③……编号。

\begin{table}
  \centering
  \begin{threeparttable}[c]
    \caption{带附注的表格示例}
    \label{tab:three-part-table}
    \begin{tabular}{ll}
      \toprule
      文件名                 & 描述                         \\
      \midrule
      thuthesis.dtx\tnote{a} & 模板的源文件，包括文档和注释 \\
      thuthesis.cls\tnote{b} & 模板文件                     \\
      thuthesis-*.bst        & BibTeX 参考文献表样式文件    \\
      \bottomrule
    \end{tabular}
    \begin{tablenotes}
      \item [a] 可以通过 xelatex 编译生成模板的使用说明文档；
        使用 xetex 编译 \file{thuthesis.ins} 时则会从 \file{.dtx} 中去除掉文档和注释，得到精简的 \file{.cls} 文件。
      \item [b] 更新模板时，一定要记得编译生成 \file{.cls} 文件，否则编译论文时载入的依然是旧版的模板。
    \end{tablenotes}
  \end{threeparttable}
\end{table}

如某个表需要转页接排，可以使用 \pkg{longtable} 宏包，需要在随后的各页上重复表的编号。
编号后跟表题（可省略）和“（续）”，置于表上方。续表均应重复表头。

\begin{longtable}{cccc}
    \caption{跨页长表格的表题}
    \label{tab:longtable} \\
    \toprule
    表头 1 & 表头 2 & 表头 3 & 表头 4 \\
    \midrule
  \endfirsthead
    \caption*{续表~\thetable\quad 跨页长表格的表题} \\
    \toprule
    表头 1 & 表头 2 & 表头 3 & 表头 4 \\
    \midrule
  \endhead
    \bottomrule
  \endfoot
  Row 1  & & & \\
  Row 2  & & & \\
  Row 3  & & & \\
  Row 4  & & & \\
  Row 5  & & & \\
  Row 6  & & & \\
  Row 7  & & & \\
  Row 8  & & & \\
  Row 9  & & & \\
  Row 10 & & & \\
\end{longtable}



\section{算法}

算法环境可以使用 \pkg{algorithms} 或者 \pkg{algorithm2e} 宏包。

\renewcommand{\algorithmicrequire}{\textbf{输入：}\unskip}
\renewcommand{\algorithmicensure}{\textbf{输出：}\unskip}

\begin{algorithm}
  \caption{Calculate $y = x^n$}
  \label{alg1}
  \small
  \begin{algorithmic}
    \REQUIRE $n \geq 0$
    \ENSURE $y = x^n$

    \STATE $y \leftarrow 1$
    \STATE $X \leftarrow x$
    \STATE $N \leftarrow n$

    \WHILE{$N \neq 0$}
      \IF{$N$ is even}
        \STATE $X \leftarrow X \times X$
        \STATE $N \leftarrow N / 2$
      \ELSE[$N$ is odd]
        \STATE $y \leftarrow y \times X$
        \STATE $N \leftarrow N - 1$
      \ENDIF
    \ENDWHILE
  \end{algorithmic}
\end{algorithm}

%%%
%% @Author           : Albert Wang
%% @Time             : 2023-01-08 21:21:02
%% @Description      :
%% @Email            : shadowofgost@outlook.com
%% @FilePath         : /Sudathesis-Bachelor-Template/contents/chap03.tex
%% @LastTime         : 2023-01-10 11:10:59
%% @LastAuthor       : Albert Wang
%% @Software         : Vscode
%% @ Copyright Notice : Copyright (c) 2023 Albert Wang 王子睿, All Rights Reserved.
%%

% !TeX root = ../main.tex

\chapter{数学符号和公式}

\section{数学符号}

中文论文的数学符号默认遵循 GB/T 3102.11—1993《物理科学和技术中使用的数学符号》
\footnote{原 GB 3102.11—1993，自 2017 年 3 月 23 日起，该标准转为推荐性标准。}。
该标准参照采纳 ISO 31-11:1992 \footnote{目前已更新为 ISO 80000-2:2019。}，
但是与 \TeX{} 默认的美国数学学会（AMS）的符号习惯有所区别。
具体地来说主要有以下差异：
\begin{enumerate}
  \item 大写希腊字母默认为斜体，如
    \begin{equation*}
      \Gamma \Delta \Theta \Lambda \Xi \Pi \Sigma \Upsilon \Phi \Psi \Omega.
    \end{equation*}
    注意有限增量符号 $\increment$ 固定使用正体，模板提供了 \cs{increment} 命令。
  \item 小于等于号和大于等于号使用倾斜的字形 $\le$、$\ge$。
  \item 积分号使用正体，比如 $\int$、$\oint$。
  \item
    偏微分符号 $\partial$ 使用正体。
  \item
    省略号 \cs{dots} 按照中文的习惯固定居中，比如
    \begin{equation*}
      1, 2, \dots, n \quad 1 + 2 + \dots + n.
    \end{equation*}
  \item
    实部 $\Re$ 和虚部 $\Im$ 的字体使用罗马体。
\end{enumerate}

以上数学符号样式的差异可以在模板中统一设置。
另外国标还有一些与 AMS 不同的符号使用习惯，需要用户在写作时进行处理：
\begin{enumerate}
  \item 数学常数和特殊函数名用正体，如
    \begin{equation*}
      \uppi = 3.14\dots; \quad
      \symup{i}^2 = -1; \quad
      \symup{e} = \lim_{n \to \infty} \left( 1 + \frac{1}{n} \right)^n.
    \end{equation*}
  \item 微分号使用正体，比如 $\dif y / \dif x$。
  \item 向量、矩阵和张量用粗斜体（\cs{symbf}），如 $\symbf{x}$、$\symbf{\Sigma}$、$\symbfsf{T}$。
  \item 自然对数用 $\ln x$ 不用 $\log x$。
\end{enumerate}


英文论文的数学符号使用 \TeX{} 默认的样式。
如果有必要，也可以通过设置 \verb|math-style| 选择数学符号样式。

关于量和单位推荐使用
\href{http://mirrors.ctan.org/macros/latex/contrib/siunitx/siunitx.pdf}{\pkg{siunitx}}
宏包，
可以方便地处理希腊字母以及数字与单位之间的空白，
比如：
\SI{6.4e6}{m}，
\SI{9}{\micro\meter}，
\si{kg.m.s^{-1}}，
\SIrange{10}{20}{\degreeCelsius}。



\section{数学公式}

数学公式可以使用 \env{equation} 和 \env{equation*} 环境。
注意数学公式的引用应前后带括号，通常使用 \cs{eqref} 命令，比如式\eqref{eq:example}。
\begin{equation}
  \frac{1}{2 \uppi \symup{i}} \int_\gamma f = \sum_{k=1}^m n(\gamma; a_k) \mathscr{R}(f; a_k).
  \label{eq:example}
\end{equation}

多行公式尽可能在“=”处对齐，推荐使用 \env{align} 环境。
\begin{align}
  a & = b + c + d + e \\
    & = f + g
\end{align}



\section{数学定理}

定理环境的格式可以使用 \pkg{amsthm} 或者 \pkg{ntheorem} 宏包配置。
用户在导言区载入这两者之一后，模板会自动配置 \env{thoerem}、\env{proof} 等环境。

\begin{theorem}[Lindeberg--Lévy 中心极限定理]
  设随机变量 $X_1, X_2, \dots, X_n$ 独立同分布， 且具有期望 $\mu$ 和有限的方差 $\sigma^2 \ne 0$，
  记 $\bar{X}_n = \frac{1}{n} \sum_{i+1}^n X_i$，则
  \begin{equation}
    \lim_{n \to \infty} P \left(\frac{\sqrt{n} \left( \bar{X}_n - \mu \right)}{\sigma} \le z \right) = \Phi(z),
  \end{equation}
  其中 $\Phi(z)$ 是标准正态分布的分布函数。
\end{theorem}
\begin{proof}
  Trivial.
\end{proof}

同时模板还提供了 \env{assumption}、\env{definition}、\env{proposition}、
\env{lemma}、\env{theorem}、\env{axiom}、\env{corollary}、\env{exercise}、
\env{example}、\env{remar}、\env{problem}、\env{conjecture} 这些相关的环境。

%%%
%% @Author           : Albert Wang
%% @Time             : 2023-04-06 22:03:52
%% @Description      :
%% @Email            : shadowofgost@outlook.com
%% @FilePath         : /sudathesis-soochow-university-latex-template/contents/chap04.tex
%% @LastTime         : 2023-05-17 16:26:57
%% @LastAuthor       : Albert Wang
%% @Software         : Vscode
%% @ Copyright Notice : Copyright (c) 2023 Albert Wang 王子睿, All Rights Reserved.
%%

% !TeX root = ../thuthesis-example.tex

\chapter{引用文献的标注}

模板支持 BibTeX 和 BibLaTeX 两种方式处理参考文献。
下文主要介绍 BibTeX 配合 \pkg{natbib} 宏包的主要使用方法。


\section{顺序编码制}

在顺序编码制下，默认的 \cs{cite} 命令同 \cs{citep} 一样，序号置于方括号中，
引文页码会放在括号外。
统一处引用的连续序号会自动用短横线连接。

\thusetup{
  cite-style = super,
}
\noindent
\begin{tabular}{l@{\quad$\Rightarrow$\quad}l}
  \verb|\cite{zhangkun1994}|               & \cite{zhangkun1994}               \\
  \verb|\citet{zhangkun1994}|              & \citet{zhangkun1994}              \\
  \verb|\citep{zhangkun1994}|              & \citep{zhangkun1994}              \\
  \verb|\cite[42]{zhangkun1994}|           & \cite[42]{zhangkun1994}           \\
  \verb|\cite{zhangkun1994,zhukezhen1973}| & \cite{zhangkun1994,zhukezhen1973} \\
\end{tabular}


也可以取消上标格式，将数字序号作为文字的一部分。
建议全文统一使用相同的格式。

\thusetup{
  cite-style = inline,
}
\noindent
\begin{tabular}{l@{\quad$\Rightarrow$\quad}l}
  \verb|\cite{zhangkun1994}|               & \cite{zhangkun1994}               \\
  \verb|\citet{zhangkun1994}|              & \citet{zhangkun1994}              \\
  \verb|\citep{zhangkun1994}|              & \citep{zhangkun1994}              \\
  \verb|\cite[42]{zhangkun1994}|           & \cite[42]{zhangkun1994}           \\
  \verb|\cite{zhangkun1994,zhukezhen1973}| & \cite{zhangkun1994,zhukezhen1973} \\
\end{tabular}



\section{著者-出版年制}

著者-出版年制下的 \cs{cite} 跟 \cs{citet} 一样。

\thusetup{
  cite-style = author-year,
}
\noindent
\begin{tabular}{@{}l@{$\Rightarrow$}l@{}}
  \verb|\cite{zhangkun1994}|                & \cite{zhangkun1994}                \\
  \verb|\citet{zhangkun1994}|               & \citet{zhangkun1994}               \\
  \verb|\citep{zhangkun1994}|               & \citep{zhangkun1994}               \\
  \verb|\cite[42]{zhangkun1994}|            & \cite[42]{zhangkun1994}            \\
  \verb|\citep{zhangkun1994,zhukezhen1973}| & \citep{zhangkun1994,zhukezhen1973} \\
\end{tabular}

\vskip 2ex
\thusetup{
  cite-style = super,
}
注意，引文参考文献的每条都要在正文中标注
\cite{zhangkun1994,zhukezhen1973,dupont1974bone,zhengkaiqing1987,%
  jiangxizhou1980,jianduju1994,merkt1995rotational,mellinger1996laser,%
  bixon1996dynamics,mahui1995,carlson1981two,taylor1983scanning,%
  taylor1981study,shimizu1983laser,atkinson1982experimental,%
  kusch1975perturbations,guangxi1993,huosini1989guwu,wangfuzhi1865songlun,%
  zhaoyaodong1998xinshidai,biaozhunhua2002tushu,chubanzhuanye2004,%
  who1970factors,peebles2001probability,baishunong1998zhiwu,%
  weinstein1974pathogenic,hanjiren1985lun,dizhi1936dizhi,%
  tushuguan1957tushuguanxue,aaas1883science,fugang2000fengsha,%
  xiaoyu2001chubanye,oclc2000about,scitor2000project%
}。

\chapter{绪论}
\section{研究背景}
\subsection{宏观政策}
近年来，我国政府积极推动新能源汽车产业的发展，将其作为战略性新兴产业加以重点支持和引导。根据中国汽车工业协会的数据，新能源汽车产量连续8年保持全球第一，在2022年达到705.8万辆，同比增长96.9\%；销量为688.7万辆，同比增长93.4\%，市场渗透率达到25.6\%。这些数据表明我国新能源汽车市场已经进入规模扩张的爆发期和全面市场化的拓展期。

这一举措是基于多方面考虑的宏观政策背景所推动的。首先，随着全球能源环境问题的日益突出，为应对气候变化、减少污染排放，我国积极响应国际社会的环保呼吁，加大了对清洁能源的投入和支持力度。新能源汽车作为清洁能源的代表，具有减少碳排放、降低能源消耗等显著优势，符合国家能源转型和环境保护的发展方向。

其次，新能源汽车产业被视为我国经济结构调整和升级的重要方向之一。通过发展新能源汽车产业，可以促进汽车产业向高端、智能化方向发展，提升我国制造业的整体竞争力。同时，新能源汽车产业链涉及到电池、充电设施等多个细分领域，有助于形成完整的产业生态体系，推动相关产业链的协同发展。

此外，新能源汽车产业也被看作是扩大内需、促进经济增长的重要引擎之一。政府通过一系列政策措施，如购车补贴、免征购置税等，鼓励消费者购买新能源汽车，推动市场需求的扩大，促进了相关产业链的发展和就业增长。

综上所述，我国政府积极推动新能源汽车产业的发展，为新能源汽车上市公司提供了良好的发展环境和政策支持。在这一背景下，研究我国新能源上市公司的最优资本结构，以比亚迪公司为例，对于指导产业健康发展、提高企业竞争力具有重要意义。

\subsection{行业发展现状}
新能源行业是中国经济发展的重要支柱，也是实现碳达峰、碳中和目标的关键领域。近年来，中国政府出台了一系列政策，大力推动新能源行业的发展。这些政策包括新能源汽车购置税减免政策、新能源发电补贴政策等，旨在通过政策引导，推动新能源技术的研发和应用，提高新能源产业的竞争力。

在这个背景下，中国的新能源行业发展迅速。据统计，2021年，全国新能源发电装机容量约占全国电源总容量的26.6\%，其中风电装机容量为$3.28 \times 10^{8} \, \text{kW}$，光伏发电为$3.06 \times 10^{8} \, \text{kW}$。全国新能源发电量为$9.785 \times 10^{11} \, \text{kW} \cdot \text{h}$，约占总发电量的11.7\%，其中风电发电量为$6.526\times 10^{11} \, \text{kW} \cdot \text{h}$，光伏发电量为$3.259 \times 10^{11} \, \text{kW} \cdot \text{h}$。这些数据充分证明了新能源行业在中国的发展势头。\cite{Dong2022}

此外，新能源汽车产业也取得了显著的发展。2022年，全国新能源汽车产量700.3万辆，比上年增长90.5\%。这一数据显示出中国新能源汽车市场的巨大潜力和发展空间。

然而，新能源行业的发展也面临一些挑战。首先，新能源技术的研发和应用还需要进一步加强，以提高新能源的效率和经济性。其次，新能源产业的发展需要大量的资金投入，而资本市场的支持是关键。因此，如何构建最优的资本结构，以满足新能源企业的发展需要，是一个重要的课题。

以比亚迪公司为例，作为我国新能源上市公司的代表，其资本结构的优化对于公司的发展具有重要意义。比亚迪公司是一家致力于“用技术创新，满足人们对美好生活的向往”的高新技术企业，其业务布局涵盖电子、汽车、新能源和轨道交通等领域。因此，研究比亚迪公司的资本结构，不仅可以为比亚迪公司的发展提供参考，也可以为我国新能源行业的发展提供借鉴。

总的来说，我国新能源行业的发展前景广阔，但也面临一些挑战。因此，我们需要进一步加强新能源技术的研发和应用，优化新能源企业的资本结构，以推动新能源行业的持续健康发展。

\subsection{企业发展需要}
资本结构优化策略是企业发展中的重中之重。比亚迪作为我国新能源产业的代表，其资本结构优化策略将成为研究的重点。需要深入探讨比亚迪在资本结构方面的选择和调整，包括债务与权益的比例、债务结构、资本成本管理等方面的策略和实践。通过对比亚迪公司的经验研究，可以为其他新能源上市公司提供优化资本结构的启示和借鉴。

另外，研究还需要关注融资渠道与方式。这包括比亚迪公司选择的融资渠道和方式，如股权融资、债券融资、银行贷款等。需要分析各种融资方式对企业财务状况、成本结构以及未来发展战略的影响，探讨如何有效利用各种融资方式来支持企业的发展需求。

此外，资本结构与企业绩效之间的关系也是研究的重要内容。通过分析比亚迪公司的财务指标、股价表现、市场份额等方面的数据，评估资本结构对企业盈利能力、成长性和市场竞争力的影响。这有助于找到最佳的资本结构配置方案，提高企业长期价值。

最后，研究还应关注风险管理与资本结构之间的关系，技术创新与资本支持的相互影响，以及战略规划与资本配置之间的协调。这些方面的研究将有助于全面理解新能源上市公司在资本结构优化方面的挑战与机遇，推动我国新能源产业的健康发展和可持续增长。
\section{研究目的与研究意义}
\subsection{研究目的}
优化新能源上市公司的资本结构是本文的核心目标。以比亚迪公司为案例，详细分析其资本结构现状，并深入剖析存在的问题。基于多目标规划法，设计资本结构优化模型，为比亚迪公司提供参考，实现资本结构优化，最终达到企业价值最大化的目标。

同时，本文旨在为我国其他新能源上市公司提供理论参考。在新能源产业普遍存在高债务融资和债务比例失衡等问题的背景下，本文以比亚迪公司为典型案例，紧密结合理论知识与新能源行业情况，提出具体的资本结构优化建议。这些建议对其他新能源公司具有借鉴意义，帮助它们找出问题并加以改正，推动整个新能源市场健康发展
\subsection{研究意义}
\subsubsection{理论研究意义}
 丰富和发展资本结构理论：通过对新能源上市公司，特别是比亚迪公司的资本结构进行深入研究，我们可以丰富和发展现有的资本结构理论。这项研究将有助于我们更好地理解新能源行业的特殊性，以及这些特殊性如何影响公司的资本结构。这不仅可以为我们提供新的理论视角，也可以为我们提供新的研究方法。

提供理论依据：这项研究将为新能源上市公司的资本结构决策提供理论依据。通过优化资本结构，新能源上市公司可以降低财务风险，提高企业价值，从而实现可持续发展。这将有助于新能源上市公司制定更科学、更合理的财务决策。

推动相关领域的研究：这项研究还将推动新能源、财务管理、企业战略等相关领域的研究。通过对新能源上市公司的资本结构进行深入研究，我们可以发现新的研究问题，开启新的研究方向。这将有助于我们更全面、更深入地理解新能源行业，也将有助于我们更好地服务新能源行业的发展。
\subsubsection{现实研究意义}
丰富和发展资本结构理论：通过对新能源上市公司，特别是比亚迪公司的资本结构进行深入研究，我们可以丰富和发展现有的资本结构理论。这项研究将有助于我们更好地理解新能源行业的特殊性，以及这些特殊性如何影响公司的资本结构。这不仅可以为我们提供新的理论视角，也可以为我们提供新的研究方法。

提供理论依据：这项研究将为新能源上市公司的资本结构决策提供理论依据。通过优化资本结构，新能源上市公司可以降低财务风险，提高企业价值，从而实现可持续发展。这将有助于新能源上市公司制定更科学、更合理的财务决策。

推动相关领域的研究：这项研究还将推动新能源、财务管理、企业战略等相关领域的研究。通过对新能源上市公司的资本结构进行深入研究，我们可以发现新的研究问题，开启新的研究方向。这将有助于我们更全面、更深入地理解新能源行业，也将有助于我们更好地服务新能源行业的发展。

\section{国内外文献综述}
\chapter{资本结构相关理论}
\section{理论概述}
\subsection{资本结构}
\subsection{资本结构优化}
\section{资本结构相关理论}
\subsection{MM理论及其修正}
\subsection{权衡理论}
\subsection{优序融资理论}
\section{资本结构优化相关理论}
\subsection{资本结构优化的目标}
\subsection{资本结构静态优化}
\subsection{资本结构动态优化}
\chapter{比亚迪概况}
\section{汽车行业概况}
\subsection{汽车行业资本结构特征}
中国上市公司汽车行业在资本结构上展现出独特的特征。首先，该行业的资本结构常常呈现较高的资产负债比例。这一现象是由汽车行业的资本密集型特性所决定的，公司需要大量的资金用于研发、生产和销售汽车产品，因此倾向于采用较高比例的债务融资来支持经营活动。

其次，中国汽车行业上市公司的资本结构通常保持相对稳健的债务结构。尽管资产负债比例较高，但许多大型企业倾向于保持一定比例的权益融资，以降低财务风险并提高偿债能力。这种债务结构的稳健性有助于企业在面对市场波动和经济不确定性时保持稳定。

此外，中国汽车行业的大型企业倾向于采取多元化融资方式。这包括银行贷款、发行债券、股权融资等多种途径，以降低融资成本、分散风险，并增强企业的融资灵活性。这种多元化的融资结构使得企业能够更好地适应市场变化和资金需求的变化。

另外，中国汽车行业上市公司的资本结构特征也受到经济周期和市场需求波动的影响较大。在行业景气期，企业可能会更加倾向于增加债务融资以支持扩张和投资；而在市场低迷时，可能会调整资本结构以降低财务风险。这种动态的资本结构调整有助于企业应对不同环境下的挑战和机遇。
\subsection{汽车行业资本结构现状}

%\subsection{资金来源}
%\subsection{债务结构}
%\subsection{股权结构}


\section{比亚迪基本情况介绍}
比亚迪公司是中国一家知名的综合性高科技企业，总部位于广东深圳。公司成立于1995年，是全球领先的新能源汽车和充电设施制造商之一，也涉足了电池、光伏发电、IT、通信等多个领域，是中国新能源汽车领域的领军企业之一。

比亚迪公司主要业务包括新能源汽车制造、电池制造、光伏发电、IT和通信等。在新能源汽车领域，比亚迪是中国最大的新能源汽车制造商之一，拥有包括电动客车、电动巴士、电动出租车、电动乘用车等在内的多款产品线。其旗下拥有比亚迪汽车股份有限公司，是中国A股市场上市的公司之一。

在电池制造方面，比亚迪是全球最大的锂电池制造商之一，产品涵盖电动车辆用电池、储能电池、手机电池等多个领域。比亚迪的锂铁磷酸铁锂电池技术在新能源汽车领域具有较高的市场份额和竞争优势。

此外，比亚迪还涉足光伏发电领域，拥有自己的光伏电池和光伏组件制造能力，并提供光伏发电系统解决方案。在IT和通信领域，比亚迪也有相关业务，包括智能手机、通信设备等产品的制造和销售。

作为一家国际化的企业，比亚迪在全球范围内建立了广泛的销售网络和合作伙伴关系，产品远销至欧洲、美洲、亚洲等多个国家和地区。公司注重科技创新和可持续发展，不断推出新产品和新技术，致力于推动清洁能源和智能交通的发展。
\subsection{经营战略}
比亚迪（BYD）作为中国一家领先的综合性高科技企业，其经营战略历程可以追溯至创立之初至今，经历了多个重要阶段。

在早期的阶段，比亚迪专注于电池制造领域，并通过不断的技术研发和产品创新，逐渐建立起自己的电池生产线。这一阶段，公司致力于提升电池技术水平和生产效率，为未来的发展奠定了坚实的基础。

随着市场的发展和技术的进步，比亚迪开始关注新能源汽车领域。2008年，公司推出了首款混合动力车型“秦”，标志着其正式进军新能源汽车市场。此后，比亚迪不断推出多款新能源汽车产品，包括电动巴士、电动出租车等，成为中国新能源汽车市场的领军企业之一。

随着国内外市场的需求增加，比亚迪开始加大国际化战略的力度。公司在全球范围内建立了广泛的销售网络和合作伙伴关系，产品远销至欧洲、美洲、亚洲等多个国家和地区。此外，比亚迪还积极开展海外投资和合作项目，进一步提升了国际影响力和竞争力。

近年来，比亚迪将科技创新作为核心驱动力，不断推出新产品和新技术。公司在电池技术、智能交通、智能制造等领域取得了一系列重要突破，推动了清洁能源和智能交通的发展。比亚迪不仅在汽车领域有所突破，还在电池、光伏发电、IT和通信等多个领域展现了强大的创新实力和竞争优势。

比亚迪经营战略历程呈现出从电池制造到新能源汽车、再到国际化扩张和科技创新的发展轨迹。公司以技术创新和市场需求为导向，不断调整战略布局，致力于成为全球领先的清洁能源和智能交通解决方案提供商。
\subsection{主营业务情况}

\subsection{内部治理结构}

\section{比亚迪资本情况介绍}

\subsection{股权结构}


\subsection{债务结构}


\subsection{长期资本}


\subsection{短期资本}

\chapter{比亚迪资本结构分析}


\section{比亚迪资本情况介绍}


\subsection{股权结构}
\begin{table}
  \centering
  \begin{threeparttable}[c]
    \caption{比亚迪2023年股权结构分布}
    \label{tab:equity-structure}
    \begin{tabular}{ccccc}
      \toprule
        股东名称 & 持股数量 & 持股比例 & 变动数量 & 变动比例 \\ 
        单位名称 & 亿元    &  \%      &  千股   &  \%      \\
      \midrule
        HKSCC NOMINEES LIMITED(H股) & 10.97  & 37.70\% & -4.437 & 0.00\% \\ 
        王传福 & 5.14  & 17.64\% & 0 & 0 \\ 
        吕向阳 & 2.39  & 8.22\% & 0 & 0 \\ 
        融捷投资控股集团有限公司 & 1.54  & 5.30\% & -971.9 & -0.63\% \\ 
        香港中央结算有限公司 & 1.04  & 3.57\% & 11548.806 & 13.30\% \\ 
        夏佐全 & 0.83  & 2.84\% & 21480.397 & 27.28\% \\ 
        王念强 & 0.18  & 0.63\% & 0 & 0 \\ 
        中央汇金资产管理有限责任公司 & 0.12  & 0.41\% & 0 & 0 \\ 
        李柯 & 0.11  & 0.37\% & -60 & -0.55\% \\ 
        中信里昂资产管理有限公司 & 0.10  & 0.35\% & 0 & 0 \\ 
      \bottomrule
    \end{tabular}
    \begin{tablenotes}
      \item [a] 数据来源：CSMAR数据库
    \end{tablenotes}
  \end{threeparttable}
\end{table}
根据表\eqref{tab:equity-structure}，我们可以发现2023年比亚迪的股权结构显示出相对稳定的特点。公司最大股东HKSCC NOMINEES LIMITED持股比例为37.70\%，虽然持股数量略有减少，但仍是最大股东。创始人王传福持股比例为17.64\%，稳定持有较大比例的股份。除此之外，其他主要股东持股比例也相对稳定，没有大幅变动。然而，部分持股数量增加较多的股东，如夏佐全和香港中央结算有限公司，显示了对公司未来发展的乐观态度。相反，持股数量减少的股东可能存在对公司风险的担忧。综合来看，比亚迪的股权结构虽有细微变动，但总体稳定，主要股东对公司的信心和预期依然较高。
\subsection{债务结构}
表\eqref{tab:current-liabilities}展示了某项目在2021年、2022年和2023年的负债情况，主要包括短期借款、交易性金融负债、应付票据、应付账款、预收款项、合同负债、应付职工薪酬、应交税费、应付利息、应付股利、其他应付款、一年内到期的非流动负债和其他流动负债等项。
\begin{table}
  \centering
  \begin{threeparttable}[c]
    \caption{比亚迪流动性负债结构}
    \label{tab:current-liabilities}
    \begin{tabular}{ccccccc}
      \toprule
      项目名称 & 2021 & ~ & 2022 & ~ & 2023 & ~ \\ 
        单位：亿元 & 数值 & 占比 & 数值 & 占比 & 数值 & 占比 \\ 
      \midrule
        短期借款 & 102.00  & 5.96\% & 51.50  & 1.55\% & 183.00  & 4.03\% \\ 
        交易性金融负债 & 0.00  & 0.00\% & 0.55  & 0.02\% & 0.08  & 0.00\% \\ 
        应付票据 & 73.30  & 4.29\% & 33.30  & 1.00\% & 40.50  & 0.89\% \\ 
        应付账款 & 732.00  & 42.81\% & 1400.00  & 42.04\% & 1940.00  & 42.73\% \\ 
        预收款项 & 0.01  & 0.00\% & 0.00  & 0.00\% & 0.00  & 0.00\% \\ 
        合同负债 & 149.00  & 8.71\% & 355.00  & 10.66\% & 347.00  & 7.64\% \\ 
        应付职工薪酬 & 58.50  & 3.42\% & 120.00  & 3.60\% & 171.00  & 3.77\% \\ 
        应交税费 & 17.80  & 1.04\% & 43.30  & 1.30\% & 78.50  & 1.73\% \\ 
        应付利息 & 0.00  & 0.00\% & 0.00  & 0.00\% & 0.00  & 0.00\% \\ 
        应付股利 & 0.00  & 0.00\% & 0.00  & 0.00\% & 0.00  & 0.00\% \\ 
        其他应付款 & 413.00  & 24.15\% & 1220.00  & 36.64\% & 1650.00  & 36.34\% \\ 
        一年内到期的非流动负债 & 130.00  & 7.60\% & 64.60  & 1.94\% & 77.40  & 1.70\% \\ 
        其他流动负债 & 37.10  & 2.17\% & 39.00  & 1.17\% & 44.50  & 0.98\% \\ 
        流动负债合计 & 1710.00  & 100.00\% & 3330.00  & 100.00\% & 4540.00  & 100.00\% \\ 
      \bottomrule
    \end{tabular}
    \begin{tablenotes}
      \item [a] 数据来源：CSMAR数据库
    \end{tablenotes}
  \end{threeparttable}
\end{table}

首先，从短期借款和交易性金融负债来看，在这三年内，短期借款有明显的波动，从2021年的102亿元增加到2023年的183亿元，而交易性金融负债则较为稳定，占比较低。

其次，应付票据和应付账款在这三年内也出现了较大的变化，应付票据从73.3亿元降低到40.5亿元，应付账款则从732亿元增加到1940亿元。这表明项目在这段时间内面临了较大的票据和账款支付压力，需要调整资金使用情况。

再者，其他应付款也出现了较大的波动，从413亿元增加到1650亿元，其他流动负债也有所增加。这可能是因为项目在这几年内有大额应付款项或其他负债产生，导致其他应付款的增加。

数据反映了该项目在这三年内负债情况的变化，需要注意短期借款、应付票据、应付账款和其他应付款等负债项的波动，合理规划资金使用和负债管理，确保项目的财务稳健和资金流动性。

\begin{table}
  \centering
  \begin{threeparttable}[c]
    \caption{比亚迪非流动性负债结构}
    \label{tab:stable-liabilities}
    \begin{tabular}{ccccccc}
      \toprule
      项目名称 & 2021 & ~ & 2022 & ~ & 2023 & ~ \\ 
        单位：亿元 & 数值 & 占比 & 数值 & 占比 & 数值 & 占比 \\ 
      \midrule
        长期借款 & 87.40  & 43.27\% & 75.90  & 19.41\% & 120.00  & 15.92\% \\ 
        应付债券 & 20.50  & 10.15\% & 0.00  & 0.00\% & 0.00  & 0.00\% \\ 
        长期负债合计 & 122.00  & 60.40\% & 102.00  & 26.09\% & 208.00  & 27.59\% \\ 
        预计负债 & 0.00  & 0.00\% & 0.00  & 0.00\% & 0.00  & 0.00\% \\ 
        递延所得税负债 & 6.10  & 3.02\% & 20.20  & 5.17\% & 39.50  & 5.24\% \\ 
        其他非流动负债 & 74.20  & 36.73\% & 269.00  & 68.80\% & 506.00  & 67.11\% \\ 
        非流动负债合计 & 202.00  & 100.00\% & 391.00  & 100.00\% & 754.00  & 100.00\% \\ 
      \bottomrule
    \end{tabular}
    \begin{tablenotes}
      \item [a] 数据来源：CSMAR数据库
    \end{tablenotes}
  \end{threeparttable}
\end{table}
从表\eqref{tab:stable-liabilities}我们发现比亚迪在过去三年的非流动负债结构中，主要涉及长期借款、应付债券、递延所得税负债和其他非流动负债等方面。首先，我们可以观察到长期借款在这三年内出现了显著的波动，从2021年的87.40亿元增加到2023年的120亿元，特别是在2022年到2023年间的变化较为明显。这可能反映了公司在资金管理和运作方面的调整和变化，可能是为了支持未来的业务发展或者优化资本结构。

另一方面，应付债券在这三年内保持为0亿元，这说明公司在这段时间内没有发行新的债券。递延所得税负债则有所增加，从2021年的6.10亿元增加到2023年的39.50亿元，这是由于公司未来所得税负债增加的预期，可能是由于盈利水平提高导致的税务压力增加。

除此之外，其他非流动负债在这三年内也出现了较大的波动，从2021年的74.20亿元增加到2023年的506亿元，增幅较大。这可能是因为公司在这段时间内承担了更多的非流动负债或者进行了大规模的业务扩张，导致负债结构发生变化。

\section{比亚迪资本结构横向分析}
\subsection{比较对象选择}
通过查阅资料，目前的新能源汽车品牌主要有蔚来、理想和小鹏。此外还有赛力斯，上汽集团等上市公司。最终我们选择长安汽车、一汽解放、上汽集团、赛力斯、蔚来、小鹏和理想作为比较的同业公司。
\subsection{同行业公司对比分析}
\begin{table}
  \centering
  \begin{threeparttable}[c]
    \caption{比亚迪与同行业优秀公司比较数据}
    \label{tab:comparing-with-other-companies}
    \begin{tabular}{ccccc}
      \toprule
        公司名称 & 资产负债率 & 流动负债比率 & 长期负债比率 & 流动比率  \\ 
      \midrule
        长安汽车 & 56.90\% & 96.21\% & 1.98\% & 127.68\% \\ 
        一汽解放 & 58.22\% & 84.30\% & 0.17\% & 125.62\% \\ 
        上汽集团 & 66.03\% & 83.20\% & 9.38\% & 107.02\% \\ 
        赛力斯 & 79.16\% & 86.81\% & 7.85\% & 83.44\% \\ 
        蔚来 & 71.28\% & 66.82\% & 15.86\% & 129.00\% \\ 
        小鹏 & 47.78\% & 66.19\% & 27.03\% & 244.74\% \\ 
        理想 & 48.37\% & 69.74\% & 21.01\% & 180.50\% \\ 
        比亚迪 & 75.42\% & 89.50\% & 2.74\% & 72.24\% \\ 
      \bottomrule
    \end{tabular}
    \begin{tablenotes}
      \item [a] 数据来源：CSMAR数据库
    \end{tablenotes}
  \end{threeparttable}
\end{table}

从表\eqref{tab:comparing-with-other-companies}我们可以发现比亚迪相对于其他汽车制造公司在财务指标上呈现出一些特点。首先，比亚迪的资产负债率相对较高，达到了75.42\%，这表明公司的资产主要依赖债务资金支持，相比之下，长安汽车和一汽解放的资产负债率较低，分别为56.90\%和58.22\%，这两家公司更多地依赖自有资金，负债较少。

其次，比亚迪的流动负债比率为89.50\%，相对较高，说明公司面临一定的流动性风险，而长安汽车和一汽解放的流动负债比率相对较低，分别为96.21\%和84.30\%，流动性风险较小。此外，比亚迪的长期负债比率为2.74\%，较低，显示出公司相对少依赖长期债务。然而，赛力斯和蔚来等公司在长期负债比率上较高，显示出对长期融资的依赖程度较大。

最后，比亚迪的流动比率为72.24\%，相对较低，表明公司的流动资产相对不足以覆盖流动负债，而蔚来、小鹏和理想的流动比率较高，显示出这些公司在短期偿付能力方面较为充裕。

综合来看，比亚迪在财务结构上与其他公司存在差异，呈现出企业整体资产负债率偏高，长期负债比率较低，但是流动比率明显低于同行业公司处于低位。
\section{比亚迪资本结构纵向分析}


\subsection{债务比率分析}
\begin{table}
  \centering
  \begin{threeparttable}[c]
    \caption{比亚迪负债比率分析}
    \label{tab:debt-rates}
    \begin{tabular}{ccccccc}
      \toprule
        比率名称 & 2018 & 2019 & 2020 & 2021 & 2022 & 2023 \\ 
      \midrule
        流动比率 & 100.43\% & 100.72\% & 106.81\% & 96.97\% & 72.24\% & 66.60\% \\ 
        速动比率 & 77.48\% & 76.64\% & 76.76\% & 71.66\% & 48.51\% & 47.27\% \\ 
        现金比率 & 9.72\% & 10.99\% & 13.15\% & 29.08\% & 15.35\% & 23.92\% \\ 
        资产负债率 & 68.81\% & 68.00\% & 67.94\% & 64.76\% & 75.42\% & 77.86\% \\
        权益乘数 & 320.58\% & 312.52\% & 311.88\% & 283.74\% & 406.84\% & 451.64\% \\ 
        产权比率 & 220.58\% & 212.52\% & 211.88\% & 183.74\% & 306.84\% & 351.64\% \\ 
        有形资产带息债务比 & 35.17\% & 41.50\% & 31.66\% & 15.58\% & 10.78\% & 15.91\% \\ 
      \bottomrule
    \end{tabular}
    \begin{tablenotes}
      \item [a] 数据来源：CSMAR数据库
    \end{tablenotes}
  \end{threeparttable}
\end{table}
通过表\eqref{tab:debt-rates}我们得出以下结果:

流动比率：这个比率显示了公司的流动资产相对于流动负债的比例。一般来说，高于1的流动比率被认为是健康的，因为这表示公司有足够的流动资产来支付其短期债务。然而，随着时间推移，比亚迪的流动比率从2018年的100.43\%下降到2023年的66.60\%，这可能表明公司在支付短期债务方面出现了挑战。

速动比率：速动比率是衡量公司快速偿付能力的指标，排除了存货这类不易快速变现的资产。和流动比率一样，速动比率也呈下降趋势，从2018年的77.48\%下降到2023年的47.27\%。这可能意味着比亚迪在快速偿付能力方面的压力逐渐增加。

现金比率：现金比率是公司现金和现金等价物占流动负债的比例，高现金比率通常被视为有利于公司应对突发情况。比亚迪的现金比率在过去几年中有波动，但总体上保持在一个相对稳定的范围内，从2018年的9.72\%上升到2021年的29.08\%，然后在2023年略有下降至23.92\%。

资产负债率：资产负债率显示了公司资产中被债务占据的比例，其下降可能暗示了公司在减少债务方面的努力。比亚迪的资产负债率在2018年至2021年之间略有下降，但在2022年和2023年开始上升，可能需要进一步关注。

权益乘数和产权比率：这两个指标都是衡量公司资本结构的指标，显示了公司资产与股东权益的关系。权益乘数和产权比率都呈现出起伏的趋势，2022年和2023年有较大幅度的增加。这可能表明公司在扩大资产规模时更多地依赖债务。

有形资产带息债务比：这个指标显示了公司有形资产与带息债务的比例，较低的比率通常被视为有利于公司财务健康。比亚迪的有形资产带息债务比在过去几年中有所下降，表明公司在控制有形资产与债务之间的关系方面取得了一定程度的成功。


比亚迪在过去几年的债务比率数据显示了一些关键趋势和挑战。公司的流动比率和速动比率逐年下降，这表明其在支付短期债务和快速偿付能力方面面临压力。然而，现金比率相对稳定，显示了一定的资金流动性。资产负债率在近年有所上升，需要关注公司在负债管理方面的策略。此外，权益乘数和产权比率的增加可能反映了公司扩张的努力，但需要注意债务水平的稳健管理以避免财务风险。值得注意的是，有形资产带息债务比在下降，显示了公司在控制有形资产与债务之间的关系方面的一定成功。综合而言，比亚迪需要继续关注债务管理，确保资金流动性和财务健康的稳定性，以支持其未来的发展和扩张战略。
\subsection{债务结构分析}
从表4.6我们可以发现比亚迪在过去六年的财务数据显示了一些关键趋势。流动负债在这段时间内有所波动，从2018年的1147.14亿元下降到2020年的1044.92亿元，然后在2021年急剧增加至1713.04亿元，继而在2022年和2023年进一步上升。相较之下，长期负债在2018年至2020年间增长，2021年开始下降，2023年稍有回升。

负债比率方面，流动负债比率从58.96\%下降至51.98\%，然后在2021年上升至57.92\%，再次在2022年和2023年上升。长期负债比率在2019年和2020年上升，随后有所下降，但在2023年略有上升。

总体来看，比亚迪的负债水平整体上呈现增长趋势，尤其是流动负债的波动较为明显。这说明企业短期偿债压力逐渐增加，需要增强现金储备应对短期债务需要。
\begin{table}
  \centering
  \begin{threeparttable}[c]
    \caption{比亚迪负债结构分析}
    \label{tab:debt-structure}
    \begin{tabular}{ccccccc}
      \toprule
        科目名称 & 2018 & 2019 & 2020 & 2021 & 2022 & 2023 \\ 
      \midrule
        流动负债（亿元） & 1147.14  & 1062.05  & 1044.92  & 1713.04  & 3333.45  & 4536.67  \\ 
        长期负债（亿元） & 139.24  & 219.16  & 236.26  & 122.05  & 102.11  & 208.22  \\ 
        总负债（亿元） & 1338.77  & 1330.40  & 1365.63  & 1915.36  & 3724.71  & 5290.86  \\ 
        总资产（亿元） & 1945.71  & 1956.42  & 2010.17  & 2957.80  & 4938.61  & 6795.48  \\ 
        流动负债比率 & 58.96\% & 54.29\% & 51.98\% & 57.92\% & 67.50\% & 66.76\% \\ 
        长期负债比率 & 7.16\% & 11.20\% & 11.75\% & 4.13\% & 2.07\% & 3.06\% \\ 
        流动负债占总负债比率 & 85.69\% & 79.83\% & 76.52\% & 89.44\% & 89.50\% & 85.75\% \\ 
        长期负债占总负债比率 & 10.40\% & 16.47\% & 17.30\% & 6.37\% & 2.74\% & 3.94\% \\ 
      \bottomrule
    \end{tabular}
    \begin{tablenotes}
      \item [a] 数据来源：CSMAR数据库
    \end{tablenotes}
  \end{threeparttable}
\end{table}

\begin{table}
  \centering
  \begin{threeparttable}[c]
    \caption{比亚迪有息债务分析}
    \label{tab:debt-with-interest}
    \begin{tabular}{ccccccc}
      \toprule
        科目名称(单位：亿元) & 2018 & 2019 & 2020 & 2021 & 2022 & 2023 \\ 
      \midrule
        短期借款 & 377.89  & 403.32  & 164.01  & 102.04  & 51.53  & 183.23  \\ 
        应付利息 & 3.90  & 5.60  & 4.14  & 0.00  & 0.00  & 0.00  \\ 
        一年内到期的非流动负债 & 74.83  & 87.47  & 114.12  & 129.83  & 64.65  & 77.40  \\ 
        长期借款 & 68.48  & 119.48  & 147.45  & 87.44  & 75.94  & 119.75  \\ 
        应付债券 & 70.77  & 99.69  & 88.80  & 20.46  & 0.00  & 0.00  \\ 
        短期有息债券 & 381.79  & 408.93  & 168.15  & 102.04  & 51.53  & 183.23  \\ 
        有息债券合计 & 595.86  & 715.56  & 518.53  & 339.78  & 192.12  & 380.39 \\ 
      \bottomrule
    \end{tabular}
    \begin{tablenotes}
      \item [a] 数据来源：CSMAR数据库
    \end{tablenotes}
  \end{threeparttable}
\end{table}

比亚迪在过去六年的债务构成呈现出一定的波动和变化。短期借款从2018年的377.89亿元起伏不定，2020年急剧下降至164.01亿元，随后在2021年和2022年有所减少，2023年略有上升至183.23亿元。与此同时，应付利息在2018年至2020年间波动较大，2021年开始为0，2022年和2023年维持在0。

一年内到期的非流动负债在过去六年中也出现了波动，2021年和2022年有所增加，2023年再次下降至77.40亿元。长期借款和应付债券则呈现出不同的趋势，长期借款在2018年至2020年间上升，2021年开始下降，2023年略有上升；应付债券在2018年至2020年间上升，然后在2021年急剧下降至20.46亿元，2022年和2023年为0。

另外，短期有息债券的情况也比较复杂，呈现出较大的波动。总体来看，有息债券合计在过去六年中总体呈现出波动趋势，2021年有较大幅度下降，2022年和2023年有所回升。

这些数据反映了比亚迪债务结构和债务规模的变化，公司可能需要更加关注不同类型债务的结构和金额分布，以确保负债管理的合理性和稳健性，从而维持财务健康和可持续发展。
\subsection{融资结构分析}
\begin{table}
  \centering
  \begin{threeparttable}[c]
    \caption{比亚迪内源融资分析}
    \label{tab:capital-from-inside}
    \begin{tabular}{ccccccc}
      \toprule
        科目名称(单位：亿元) & 2018 & 2019 & 2020 & 2021 & 2022 & 2023 \\ 
      \midrule
        盈余公积 & 38.42  & 40.99  & 44.48  & 50.09  & 68.39  & 73.74  \\ 
        未分配利润 & 204.98  & 210.56  & 244.57  & 264.56  & 409.43  & 671.24  \\ 
        留存收益 & 243.41  & 251.56  & 289.05  & 314.65  & 477.82  & 744.98  \\ 
        折旧摊销 & 94.25  & 96.24  & 123.46  & 138.33  & 197.83  & 421.60  \\ 
        内源融资合计 & 337.65  & 347.79  & 412.51  & 452.98  & 675.65  & 1166.58  \\ 
        资产总额 & 1945.71  & 1956.42  & 2010.17  & 2957.80  & 4938.61  & 6795.48  \\ 
        内源融资占比 & 17.35\% & 17.78\% & 20.52\% & 15.31\% & 13.68\% & 17.17\% \\ 
      \bottomrule
    \end{tabular}
    \begin{tablenotes}
      \item [a] 数据来源：CSMAR数据库
    \end{tablenotes}
  \end{threeparttable}
\end{table}
从表\eqref{tab:capital-from-inside}可以看出比亚迪在过去六年内部源融资（包括盈余公积、未分配利润、留存收益和折旧摊销）呈现出稳步增长的趋势。盈余公积从2018年的38.42亿元增长至2023年的73.74亿元，未分配利润从204.98亿元增长至671.24亿元，留存收益和折旧摊销也呈现出持续增长的态势。

内源融资合计从2018年的337.65亿元增长至2023年的1166.58亿元，占比方面，内源融资占比在这段时间内波动，但总体保持在相对稳定的水平，从17.35\%上升至2021年的20.52\%，然后再次下降至2023年的17.17\%。

这些数据显示了比亚迪通过内部渠道获得资金的能力逐渐增强，特别是未分配利润和留存收益等方面的增长较为显著。公司内部源融资占比的波动可能受到经营策略和财务规划的影响.从占比来说内源融资规模不断增大，速度不断增快，整体占比较为平稳。
\begin{table}
  \centering
  \begin{threeparttable}[c]
    \caption{比亚迪外源融资分析}
    \label{tab:capital-from-outside}
    \begin{tabular}{ccccccc}
      \toprule
        科目名称(单位：亿元) & 2018 & 2019 & 2020 & 2021 & 2022 & 2023 \\ 
      \midrule
        吸收权益性投资收到的现金 & 0.11  & 0.02  & 28.00  & 373.14  & 5.08  & 0.98  \\ 
        发行债券收到的现金 & 141.00  & 200.00  & 20.00  & 0.00  & 0.00  & 0.00  \\ 
        取得借款收到的现金 & 529.65  & 584.78  & 406.34  & 328.72  & 276.36  & 453.04  \\ 
        筹资活动现金流入小计 & 670.76  & 797.96  & 454.34  & 701.86  & 311.75  & 454.30  \\ 
        资产总额 & 1945.71  & 1956.42  & 2010.17  & 2957.80  & 4938.61  & 6795.48  \\ 
        外源融资占比 & 34.47\% & 40.79\% & 22.60\% & 23.73\% & 6.31\% & 6.69\% \\ 
      \bottomrule
    \end{tabular}
    \begin{tablenotes}
      \item [a] 数据来源：CSMAR数据库
    \end{tablenotes}
  \end{threeparttable}
\end{table}

比亚迪在过去六年的外源融资方面表现出一定的波动和变化。吸收权益性投资收到的现金在2018年至2020年间波动较大，2021年达到了373.14亿元的高点，然后在2022年和2023年有所下降。发行债券收到的现金在2018年和2019年较高，2020年急剧下降至20亿元，后续年份为0。取得借款收到的现金在过去六年中有所波动，2023年略有增加至453.04亿元。

筹资活动现金流入小计在2018年至2021年间波动较大，2022年和2023年有所下降。外源融资占比从2018年的34.47\%上升至2019年的40.79\%，然后在后续年份有所波动，2022年和2023年稳定在较低的水平，约为6.31\%至6.69\%。

这些数据显示了比亚迪在过去六年内外部融资的波动情况，尤其是在吸收权益性投资收到的现金和发行债券收到的现金方面的大幅波动。公司可能需要进一步关注外源融资的稳定性和可持续性，同时密切关注各种融资方式对公司财务结构和成本的影响。从中我们也可以发现外源融资规模波动较大，并且呈现整体缩减的态势。占资产总额占比在不断的缩小。

结合外源融资和内源融资来看，比亚迪逐渐从原来的外源融资逐渐转向内源融资，内源融资规模不断扩大。
\section{比亚迪资本结构存在的问题及原因}

\subsection{资产负债率过高}
近年来，比亚迪资产负债率不断上升，从2018年的68.81\%一直上升到2023年的77.86\%，高于同行业企业的资产负债率。只低于赛力斯的资产负债率，无论是和传统汽车厂商还是和新兴的新能源企业相比，资产负债率均高过10\%。这说明比亚迪和同行业公司相比面临着更高的财务风险，可能在经营中出现现金流紧张等问题，严重情况会导致资金链断裂，企业破产。
\subsection{短期债务占比大}
从占比来说，流动负债占总负债比率高达85.75\%，而长期负债占总负债比率从2020年的17.3\%降低到2023年的3.94\%。负债比率严重失衡。考虑到疫情期间经济下行压力较大，使用短期负债融资具有资金成本低廉，融资迅速的特点。但是考虑到新能源企业最近竞争激烈，不断有企业通过价格战的方式抢占市场，企业必须在未来调整负债结构，减少短期债务占总负债的比重，从而获得更大的收益。

当企业借入短期负债时，尽管能够迅速获取资金，但这也带来了巨大的偿债压力。如果企业无法及时偿还这些债务，就会面临严重的信用风险，甚至可能引发金融危机。一旦陷入“借新还旧”的局面，企业就会陷入无休止的资金循环，不得不不断借新债来还旧债，最终导致资金链断裂。这将对企业声誉造成严重损害，影响销售回款，同时也会让投资者对企业未来的发展前景感到担忧，进而使企业融资更加困难。因此，过高的短期负债比例会增加企业破产的风险，不利于比亚迪集团长期稳健的发展。
\subsection{融资渠道不合理}
从比亚迪融资来源分析，我们可以发现比亚迪融资来源逐渐从原来的外源融资转向内源融资。此外在融资渠道选择上，在2022年和2023年没有大规模的通过发行债券和股票进行融资。相比较而言，比亚迪每年都会通过银行借款进行融资，融资渠道较为单一，融资渠道并不合理。 
%\subsection{股权结构较分散}



\chapter{比亚迪资本结构存在的问题和原因分析}
在第四章中我们通过横纵向分析了比亚迪的资本结构现状。在这一章中我们将会根据上一章数据和分析总结出比亚迪资本结构存在的问题以及对存在的问题进行进一步分析，分析问题的出现的原因。 
\section{比亚迪资本结构存在的问题}

\subsection{资产负债率过高}
近年来，比亚迪资产负债率不断上升，从2018年的68.81\%一直上升到2023年的77.86\%，高于大部分同行业企业的资产负债率，只低于赛力斯的资产负债率。无论是和传统汽车厂商还是和新兴的新能源企业相比，资产负债率均高过10\%。这说明比亚迪和同行业公司相比面临着更高的财务风险，可能在经营中出现现金流紧张等问题，严重情况会导致资金链断裂，企业破产。

比亚迪作为中国目前规模最大的新能源汽车企业，月销量处于同行业第一的位置，资金需求量大，融资规模较高；同时他也涉足多个产业，比如手机和电池制造等产业，面临较大的营业风险。另一方面，作为国内大规模，市场占有率名列前茅的新能源企业，较大的企业规模是得企业在债务融资方面可以获得比股权融资更低的成本，因此企业的资产负债率也居高不下。
\subsection{短期债务占比大}
从占比来说，流动负债占总负债比率高达85.75\%，而长期负债占总负债比率从2020年的17.3\%降低到2023年的3.94\%。负债比率严重失衡。考虑到疫情期间经济下行压力较大，使用短期负债融资具有资金成本低廉，融资迅速的特点。但是考虑到新能源企业最近竞争激烈，不断有企业通过价格战的方式抢占市场，企业必须在未来调整负债结构，减少短期债务占总负债的比重，从而获得更大的收益。

当企业借入短期负债时，尽管能够迅速获取资金，但这也带来了巨大的偿债压力。如果企业无法及时偿还这些债务，就会面临严重的信用风险，甚至可能引发金融危机。一旦陷入“借新还旧”的局面，企业就会陷入无休止的资金循环，不得不不断借新债来还旧债，最终导致资金链断裂。这将对企业声誉造成严重损害，影响销售回款，同时也会让投资者对企业未来的发展前景感到担忧，进而使企业融资更加困难。因此，过高的短期负债比例会增加企业破产的风险，不利于比亚迪集团长期稳健的发展。
\subsection{融资渠道不合理}
从比亚迪融资来源分析，我们可以发现比亚迪融资来源逐渐从原来的外源融资转向内源融资，外源融资占比逐渐下降。此外在融资渠道选择上，在2022年和2023年没有大规模的通过发行债券和股票进行融资。相比较而言，比亚迪每年都会通过银行借款进行融资，融资渠道较为单一，并没有通过资本市场进行融资，融资渠道并不合理。这使得企业的负债将会受到利率风险的影响，将会一定程度上影响企业的资产负债率和企业偿债能力，进而影响企业的才稳健性。
%\subsection{股权结构较分散}

\section{存在的问题原因分析}
\subsection{短期债务对业务增长的推动更明显}
短期债务主要用于资金周转，对业务发展促进作用十分明显。2022年比亚迪短期债务迅速增长，比亚迪营业收入同比增长96\%，新能源汽车销售量同比增长208\%。2023年比亚迪短期债务迅速增长，增速慢于2022年，但是营业收入同比增长依然高达42\%，新能源汽车销售量同比增长高达62\%。结合之前的资本结构分析，我们不难得出结论，短期债务对业务的发展促进效果十分明显。市场环境上，目前新能源汽车行业竞争十分激烈，不同的新能源汽车公司通过推行不同的产品，降低企业价格来获得更大的市场份额。为了应对竞争越来越激烈，市场变化越来越迅速，比亚迪因此大量使用短期债务进行融资。
\subsection{未科学评估融资需求}
从比亚迪的行业性质，发展情况以及财务数据呈现出的以短期负债为主的融资结构我们可以发现比亚迪没有根据企业需要建立与企业需求期限相匹配的融资结构，没有科学评估企业的融资需求。在最近两年，比亚迪产品销售，特别是新能源汽车的销售出现大幅增长，产生的成本支出和劳务支出出现大幅增长。但是比亚迪新能源汽车属于制造业，投入周期长，资金需求量大，应该匹配长期投资的资金和长期债务资金。而目前大量使用短期债务融资的方式使得资金来源期限和资金来源期限不匹配，短期偿债压力大，有很高的财务风险。 
\subsection{资本市场的不完善使直接融资受限}
中国资本市场的发展尚未完善，这导致企业在融资方面受到了一定的限制。相较于发达国家的资本市场，中国的资本市场在市场规模、市场流动性、市场透明度以及金融产品创新等方面存在一定差距。这种不完善的资本市场环境使得企业在融资过程中面临诸多挑战，包括融资渠道有限、融资成本较高、融资周期较长等问题。企业在市场上融资过程中存在着融资手续繁琐，融资需要的时间长。此外，市场调剂机制效果欠佳。市场的资金的流向不能很好的与企业资金需求相匹配。对于比亚迪这种新能源制造业企业而言，需要资金投入量大，投资时间长，需要进行长期融资。但是受限于目前市场情况，企业融资以短期债务融资为主，不利于企业长远发展。

\chapter{比亚迪最优资本结构}
在上一章中我们通过横纵向分析了比亚迪的资本结构现状。在这一章中我们将会采取措施优化比亚迪的资本结构。首先我们会设定优化资本结构的原则和目标。在这个基础上我们将会搭建我们的资本结构优化模型，并使用比亚迪的财报数据计算最优的资本结构。最终根据我们的计算的结果，为比亚迪资本结构优化提供建议和具体举措。
\section{最优资本结构设定原则和目标}
\subsection{最优资本结构优化原则}
风险收益均衡是优化资本结构的重要原则之一。公司需要在追求高收益的同时，合理控制风险水平。这意味着在选择资本结构时要考虑到债务和股权的比例，以及债务的成本和偿还期限。通过平衡债务和股权的结构，公司可以降低资本成本，提高资金利用效率，但同时也需要注意债务对财务风险的影响，确保债务水平不会过高，导致偿还能力下降或财务稳健性受到威胁。

资本结构区间化。公司在资本结构设计上应该考虑到不同经济环境和行业周期的影响，因此需要建立灵活的资本结构区间。这意味着公司应该在不同的经济周期和行业变化中，根据具体情况调整资本结构。例如，在经济增长期间，公司可以适度增加债务比例，以扩大业务规模和投资项目。而在经济低迷或不确定时期，公司可能会减少债务比例，增加内部融资，以降低财务风险。

资本结构优化可操作性。公司在设计资本结构时要考虑到可操作性和执行性。这包括确保资本结构设计合理、灵活，并能够根据公司发展需求和市场环境变化进行调整。公司需要建立有效的财务规划和风险管理机制，监控资本结构的变化和影响，并制定相应的调整策略。此外，公司还应该加强内部沟通和协作，确保各部门之间的配合和执行力度，以实现资本结构优化的目标。
\subsection{最优资本结构优化目标}
企业的资本结构会受到宏观经济发展水平的影响，自身的发展情况以及企业未来长远战略目标的改变，因此，企业的最优资本结构处于不断变动之中。学术界普遍认为，公司资本结构以利润最大化、每股收益最大化和企业价值最大化为三个主要目标。尽管利润最大化以公司总体利益为出发点，但其局限性在于过于偏重短期利益，忽略了企业长期利益，也未考虑时间价值因素的影响。每股收益最大化目标则未考虑企业经营风险，企业长远收益。可能为了企业短期收益而牺牲企业的长期效益。

相比之下，企业价值最大化目标对资金的时间价值和其他相关风险因素进行了全面考量，有助于提升企业自身价值并维护长远发展。因此，新能源企业要实现持续健康发展，必须不断调整和完善自身的资本结构，也与本文研究的宗旨一致。

\section{比亚迪最优资本结构模型的建立}
\subsection{静态优化模型}
本文的静态最优资本结构模型来自于孙子婷的研究\cite{Sun2020}。

该模型满足以下假设：
\begin{enumerate}[label=(\arabic*)]
    \item 企业已公开上市，且资本结构中仅含有借入资本和自有资本；
    \item 企业借入资本是企业债权的市值；自有资本是上市的普通股票市值；
    \item 资本市场不受限制；
    \item 股票的市盈率与负债率无关。
\end{enumerate}
该模型主要表述为：
\eqref{eq:refer}。
\begin{equation}
\begin{aligned}
& \operatorname{Ln}(E)=\frac{F+I}{E}+\frac{E+I+F}{V} * \frac{100 \mathrm{~B}}{1-\mathrm{B}} 
  \label{eq:refer}
\end{aligned}
\end{equation}
\begin{equation}
\begin{aligned}
& \mathrm{F}=\mathrm{Z}+\mathrm{iL}
\label{eq:else}
\end{aligned}
\end{equation}
将公式\eqref{eq:else}带入到公式\eqref{eq:refer}中得到：
\begin{equation}
\begin{aligned}
& \operatorname{Ln}(E)=\frac{Z+iL+I}{E}+\frac{E+I+Z+iL}{V} * \frac{100 \mathrm{~B}}{1-\mathrm{B}} 
  \label{eq:main}
\end{aligned}
\end{equation}
其中，$E$为税前利润总额，$I$为企业利息，计算时以财务费用代替；$V$为公司总资产；$F$为公司固定成本，固定资产折旧及长期债务利息；$Z$为资产折旧及摊销；$L$为长期债务总额，计算时用非流动负债替代；$i$为长期债务利息率；$B$为企业最优资产负债率。
%\section{最优资本结构}
根据上述模型，结合比亚迪财报数据内容，计算得出表\eqref{tab:best-structure}中的变量数据。长期负债利率以中国银行间同业拆借中心的长期利率为基础估算。 
\begin{table}
  \centering
  \begin{threeparttable}[c]
    \caption{静态区间相关计算数据}
    \label{tab:best-structure}
    \begin{tabular}{ccccccc}
      \toprule
        科目名称(单位：亿元) & 2018 & 2019 & 2020 & 2021 & 2022 & 2023 \\ 
      \midrule
        税前利润总额E & 43.86  & 24.31  & 68.83  & 45.18  & 210.80  & 372.69  \\ 
        企业利息I & 31.19  & 34.87  & 31.24  & 19.08  & 13.16  & 18.28  \\ 
        企业总资产V & 1945.71  & 1956.42  & 2010.17  & 2957.80  & 4938.61  & 6795.48  \\ 
        资产折旧和摊销Z & 94.25  & 96.24  & 123.46  & 138.33  & 197.83  & 421.60  \\ 
        长期负债L & 139.24  & 219.16  & 236.26  & 122.05  & 102.11  & 208.22  \\ 
        长期负债利率i & 4.50\% & 4.00\% & 3.75\% & 3.75\% & 3.50\% & 3.50\% \\ 
      \bottomrule
    \end{tabular}
    \begin{tablenotes}
      \item [a] 数据来源：CSMAR数据库
    \end{tablenotes}
  \end{threeparttable}
\end{table}
根据公式\eqref{eq:main}计算可得，通过最优资本结构模型，我们可以计算出2018-2023年比亚迪公司所对应的最优资本结构，如下表\eqref{tab:best-structure-result}所示：
\begin{table}
  \centering
  \begin{threeparttable}[c]
    \caption{静态区间相关计算数据}
    \label{tab:best-structure-result}
    \begin{tabular}{ccccccc}
      \toprule
        科目名称 & 2018 & 2019 & 2020 & 2021 & 2022 & 2023 \\ 
      \midrule
        行业资产负债率 & 52.85\% & 53.51\% & 56.65\% & 56.83\% & 60.40\% & 64.35\% \\ 
        比亚迪原始资产负债率 & 68.81\% & 68.00\% & 67.94\% & 64.76\% & 75.42\% & 77.86\% \\ 
        比亚迪最优资产负债率 & 68.03\% & 65.39\% & 63.69\% & 72.69\% & 72.54\% & 65.73\% \\ 
      \bottomrule
    \end{tabular}
    \begin{tablenotes}
      \item [a] 数据来源：CSMAR数据库
    \end{tablenotes}
  \end{threeparttable}
\end{table}

根据表\eqref{tab:best-structure-result}中的结果显示，除了2021年比亚迪资产负债率低于最优的资产负债率，其他时间比亚迪的资产负债率均高于最优资产负债率，尤其是2023年，比亚迪的资产负债率比最优资产负债率高12\%。并且比亚迪的最优资产负债率均高于同行业企业加权平均后的行业资产负债率。这说明比亚迪需要在未来降低企业的资产负债率，过高的资产负债率加剧了企业正常的运营风险。
\subsection{动态优化模型}
资本结构的动态优化是在静态优化的基础上进行调整，使得一个原本固定的静态优化结果变化成一个动态优化区间。本文将通过熵权法确定相关影响因素对于资本结构优化的影响权重，从而判断出关键的影响因素，确定调整值，得到更合理的动态调整区间。
\begin{table}
  \centering
  \begin{threeparttable}[c]
    \caption{指标层级}
    \label{tab:indicator-hierarchy}
    \begin{tabular}{ccc}
      \toprule
        一级指标 & 二级指标 & 三级指标 \\ 
      \midrule
        \multirow{9}{*}{企业价值最大化} & 外部因素 & GDP \\ 
                                    & 企业规模 & 企业总资产 \\ 
                        & \multirow{2}{*}{盈利能力} & 净资产收益率 \\ 
                        &                          & 总资产收益率 \\
                        & \multirow{2}{*}{成长性} & 营业收入增长率 \\ 
                        &                         & 净利润增长率 \\
                        & \multirow{2}{*}{营运能力} & 总资产周转率 \\ 
                        &                          & 应收账款周转率 \\
      \bottomrule
    \end{tabular}
  \end{threeparttable}
\end{table}

如表\eqref{tab:indicator-hierarchy}所示，我们首先构建了一个三层指标模型，这样可以更方便地计算权重指标，为确定主要影响因素提供依据。

在第一层指标中，我们选择了企业价值最大化作为最主要的指标。其次，在第二层指标中，我们考虑了外部因素和内部企业的盈利能力以及成长性。最后，在第三层指标的选择上，我们选取了GDP数据作为外部因素的计算指标，同时还包括了企业总资产、资产收益率等内部数据作为内部定量计算指标。通过这些指标的选取，我们可以更轻松地计算出动态指标数据。

在确定指标后，我们还需要对各个指标的正负性进行评估，因为熵权法的计算需要根据指标的正负性进行不同的数据处理。在这里，我们通过分析发现三级指标比如净资产收益率，营业收入增长率等指标对一级指标企业价值最大化都有正向作用，因此我们所采用的三级指标均属于正向指标。

通过这样的指标选择和正负性判断，我们可以更加准确地计算出各项指标的权重，为企业资本结构的动态优化提供可靠的数据支持和决策依据。

\begin{table}
  \centering
  \begin{threeparttable}[c]
    \caption{比亚迪公司2018-2023年相关动态指标数据}
    \label{tab:indicator-hierarchy-initial-data}
    \begin{tabular}{ccccccc}
      \toprule
        指标名称 & 2018 & 2019 & 2020 & 2021 & 2022 & 2023 \\ 
      \midrule
        国内生产总值（千亿元） & 919.28  & 986.52  & 1013.57  & 1143.67  & 1210.21  & 1260.58  \\ 
        企业总资产（亿元） & 1945.71  & 1956.42  & 2010.17  & 2957.80  & 4938.61  & 6795.48  \\ 
        净资产收益率 & 0.0586  & 0.0338  & 0.0933  & 0.0381  & 0.1459  & 0.2083  \\ 
        总资产收益率 & 0.0183  & 0.0108  & 0.0299  & 0.0134  & 0.0359  & 0.0461  \\ 
        营业收入增长率 & 0.1792  & 0.0720  & 0.1585  & 0.3065  & 0.3356  & 0.1103  \\ 
        净利润增长率 & 0.1087  & -0.0914  & -0.5584  & -0.5486  & 0.2757  & -0.1717  \\ 
        总资产周转率 & 0.6980  & 0.6547  & 0.7896  & 0.8701  & 1.0741  & 1.0266  \\ 
        应收账款周转率 & 2.5712  & 2.7407  & 3.6782  & 5.5802  & 11.2963  & 11.9632 \\ 
      \bottomrule
    \end{tabular}
    \begin{tablenotes}
      \item [a] 数据来源：CSMAR数据库
    \end{tablenotes}
  \end{threeparttable}
\end{table}
如表\eqref{tab:indicator-hierarchy-initial-data}所示，我们获取了相关指标数据以便下一步进行计算。

\subsubsection{数据标准化处理}
根据熵权法原理，给定了n 个样本， m 个指标：$\left\{ X_{1},X_{2},\cdot\cdot\cdot,X_{m}\right\}$ ，形成原始数据矩阵如下，$x_{ij}$ 为第 i 个样本的第 j 个指标的数值：
\begin{equation}
\begin{aligned}
X=\begin{pmatrix} x_{11} & x_{12} & \cdots & x_{1m} \\ x_{21} & x_{22} & \cdots & x_{2m} \\ \vdots & \vdots & \ddots & \vdots \\ x_{31}& x_{32} & \cdots &x_{nm} \\ \end{pmatrix}\quad (x_{ij},\quad i=1,2,...,n\quad j=1,2,...,m)
  \label{eq:matrix-x}
\end{aligned}
\end{equation}

在本文的模型中，我们共选取了2018-2023年的数据，也就是共6个样本数据，8个指标。

将各个指标进行标准化处理，不同类型的指标，需要按照不同的方式分别进行标准化:
\begin{equation}
\begin{aligned}
y_{ij}= \begin{cases} \frac{x_{ij}-min(X_{j})}{max(X_{j})-min(X_{j})},\quad if \quad X_{j}$为正向指标$\\ \frac{max(X_{j})-x_{ij}}{max(X_{j})-min(X_{m})},\quad if \quad X_{j}$为负向指标, $\end{cases} 
  \label{eq:standard}
\end{aligned}
\end{equation}
由上文所知，所有指标均是正向指标，所以我们以正向指标公式进行计算。
原始指标$\left\{ X_{1},X_{2},\cdot\cdot\cdot,X_{m}\right\}$转化为标准化指标 $\left\{ Y_{1},Y_{2},\cdot\cdot\cdot,Y_{m}\right\}$，$y_{ij}$ 为标准化后的第$i$个样本的第$j$个指标的数值：
\begin{equation}
\begin{aligned}
X=\begin{pmatrix} x_{11} & x_{12} & \cdots & x_{1m} \\ x_{21} & x_{22} & \cdots & x_{2m} \\ \vdots & \vdots & \ddots & \vdots \\ x_{31}& x_{32} & \cdots &x_{nm} \\ \end{pmatrix}\quad  \Rightarrow \quad Y= \begin{pmatrix} y_{11} & y_{12} & \cdots & y_{1m} \\ y_{21} & y_{22} & \cdots & y_{2m} \\ \vdots & \vdots & \ddots & \vdots \\ y_{31}& y_{32} & \cdots &y_{nm} \\ \end{pmatrix}
  \label{eq:matrix}
\end{aligned}
\end{equation}
以上面的公式为基础，基于表\eqref{tab:indicator-hierarchy-initial-data}的数据进行计算，我们得到了标准化后的数据。 
\begin{table}
  \centering
  \begin{threeparttable}[c]
    \caption{比亚迪公司2018-2023年标准化动态指标数据}
    \label{tab:indicator-hierarchy-standard-data}
    \begin{tabular}{ccccccc}
      \toprule
        指标名称 & 2018 & 2019 & 2020 & 2021 & 2022 & 2023 \\ 
      \midrule
        国内生产总值（千亿元） & 0.0001  & 0.1970  & 0.2763  & 0.6575  & 0.8524  & 1.0000  \\ 
        企业总资产（亿元） & 0.0001  & 0.0427  & 0.2837  & 0.4335  & 0.5436  & 1.0000  \\ 
        净资产收益率（\%） & 0.4051  & 0.0001  & 0.6494  & 0.6656  & 0.9776  & 1.0000  \\ 
        总资产收益率（\%） & 0.4350  & 0.0001  & 0.6778  & 0.6706  & 1.0000  & 0.9618  \\ 
        营业收入增长率（\%） & 0.2242  & 1.0000  & 0.3571  & 0.2240  & 0.3600  & 0.0001  \\ 
        净利润增长率（\%） & 0.8325  & 0.0001  & 0.7612  & 0.9045  & 1.0000  & 0.9050  \\ 
        总资产周转率（\%） & 0.0001  & 0.3162  & 0.4653  & 0.6623  & 0.8372  & 1.0000  \\ 
        应收账款周转率（\%） & 0.0001  & 1.0000  & 0.7197  & 0.6425  & 0.4583  & 0.6480 \\ 
      \bottomrule
    \end{tabular}
  \end{threeparttable}
\end{table}

由表\eqref{tab:indicator-hierarchy-standard-data}所示，我们得到了经过标准化计算的数据。

\subsubsection{计算信息熵与权重}
首先，我们需要计算某项指标$X_j$下第$i$个样本占该指标的比重，$p_{ij}$为第$i$个样本在$j$指标当中的比重值：
\begin{equation}
\begin{aligned}
p_{ij}=\frac{y_{ij}}{\sum_i^n y_{ij} } ,\quad i=1,2,...,n;\quad j=1,2,...,m
  \label{eq:pij}
\end{aligned}
\end{equation}

在完成了$p_{ij}$的计算后，我们需要进一步计算某项指标$X_j$的熵值：
\begin{equation}
\begin{aligned}
E_j=-\frac{1}{ln(n)}\sum_{i=1}^np_{ij}ln(p_{ij})
  \label{eq:ej}
\end{aligned}
\end{equation}

其中， n 为样本数量，一般而言， $0\leq E_j\leq1$

根据上述信息熵的计算公式，计算出各个指标的信息熵为$E_1,E_2,…,E_m$后，我们将用其来计算各个指标$X_j$的权重
\begin{equation}
\begin{aligned}
W_j=\frac{1-E_j}{m-\sum E_j }\quad (j=1,2,...,m)
  \label{eq:wj}
\end{aligned}
\end{equation}

\begin{table}
  \centering
  \begin{threeparttable}[c]
    \caption{比亚迪公司信息熵指标}
    \label{tab:indicator-hierarchy-result-data}
    \begin{tabular}{cccc}
      \toprule
        指标名称 & 信息熵e & 信息效用d & 信息权重w \\ 
      \midrule
        国内生产总值（千亿元） & 0.8139  & 0.1861  & 8.55\% \\ 
        企业总资产（亿元） & 0.5521  & 0.4479  & 20.59\% \\ 
        净资产收益率（\%） & 0.6919  & 0.3081  & 14.16\% \\ 
        总资产收益率（\%） & 0.7608  & 0.2392  & 10.99\% \\ 
        营业收入增长率（\%） & 0.7939  & 0.2061  & 9.47\% \\ 
        净利润增长率（\%） & 0.7617  & 0.2383  & 10.95\% \\ 
        总资产周转率（\%） & 0.7872  & 0.2128  & 9.78\% \\ 
        应收账款周转率（\%） & 0.6629  & 0.3371  & 15.49\% \\ 
      \bottomrule
    \end{tabular}
  \end{threeparttable}
\end{table}

根据表\eqref{tab:indicator-hierarchy-standard-data}的数据，结合公式\eqref{eq:pij},公式\eqref{eq:ej}和公式\eqref{eq:wj}，进行计算，并对计算结果进行非负化处理后，我们得到了表\eqref{tab:indicator-hierarchy-result-data}的结果
\begin{table}
  \centering
  \begin{threeparttable}[c]
    \caption{指标层级权重}
    \label{tab:indicator-hierarchy-result}
    \begin{tabular}{ccccc}
      \toprule
        一级指标 & 二级指标 & 权重 & 三级指标 & 权重 \\ 
      \midrule
        \multirow{9}{*}{企业价值最大化} & 外部因素 &8.55\% & GDP  & 8.55\% \\ 
        & 企业规模 & 20.59\% & 企业总资产& 20.59\% \\
        & \multirow{2}{*}{盈利能力}&\multirow{2}{*}{25.15\%} & 净资产收益率 &14.16\% \\ 
        &                         &  & 总资产收益率 & 10.99\% \\ 
        & \multirow{2}{*}{成长性} &\multirow{2}{*}{20.42\%}&营业收入增长率&9.47\% \\ 
        &                         &   & 净利润增长率& 10.95\% \\ 
        & \multirow{2}{*}{营运能力}&\multirow{2}{*}{25.27\%}& 总资产周转率& 9.78\% \\
        &                         &    & 应收账款周转率& 15.49\% \\ 
      \bottomrule
    \end{tabular}
  \end{threeparttable}
\end{table}

根据表\eqref{tab:indicator-hierarchy-result-data}的结果，我们便可以得到各级指标对比亚迪公司的影响权重，如上述表\eqref{tab:indicator-hierarchy-result}所示。

对表\eqref{tab:indicator-hierarchy-result}中的数据进行分析我们可以发现，宏观影响因素对于企业价值最大化的影响最小，然而，微观因素的四个因素均对企业价值最大化这一目标有着较大的影响。因此，在以企业价值最大调整比亚迪公司资本结构优化时，我们需要更多的从盈利能力和营运能力两个角度去考虑。宏观环境问题固然影响很小，但是我们无法改变宏观环境，所以我们需要针对性地去根据宏观环境调整我们的优化措施。 
\subsubsection{确定最优资本结构区间}
在我们之前使用的熵权法的基础上，我们选取三家行业上市公司参与计算，求得三家公司的信息熵、信息效用和信息权重的均值。并选择于比亚迪的信息效用对比，从而得出比亚迪的最优资本结构区间。
\begin{table}
  \centering
  \begin{threeparttable}[c]
    \caption{三家公司信息熵指标}
    \label{tab:indicator-hierarchy-multiple-data}
    \begin{tabular}{cccc}
      \toprule
        指标名称 & 信息熵e & 信息效用d & 信息权重w \\ 
      \midrule
        国内生产总值（千亿元） & 0.8173  & 0.1827  & 10.73\% \\ 
        企业总资产（亿元） & 0.8026  & 0.1974  & 11.59\% \\ 
        净资产收益率（\%） & 0.7899  & 0.2101  & 12.34\% \\ 
        总资产收益率（\%） & 0.7918  & 0.2082  & 12.23\% \\ 
        营业收入增长率（\%） & 0.8416  & 0.1584  & 9.30\% \\ 
        净利润增长率（\%） & 0.6361  & 0.3639  & 21.36\% \\
        总资产周转率（\%） & 0.8007  & 0.1993  & 11.70\% \\ 
        应收账款周转率（\%） & 0.8170  & 0.1830  & 10.75\% \\ 
      \bottomrule
    \end{tabular}
  \end{threeparttable}
\end{table}

\begin{table}
  \centering
  \begin{threeparttable}[c]
    \caption{指标差额分析}
    \label{tab:indicator-hierarchy-different-data}
    \begin{tabular}{ccccc}
      \toprule
        指标名称 & 比亚迪信息效用 & 三家企业信息效用 & 信息权重 & 调整额度 \\ 
      \midrule
        国内生产总值（千亿元） & 0.1861  & 0.1827  & 10.73\% & 0.0004  \\ 
        企业总资产（亿元） & 0.4479  & 0.1974  & 11.59\% & 0.0290  \\ 
        净资产收益率（\%） & 0.3081  & 0.2101  & 12.34\% & 0.0121  \\ 
        总资产收益率（\%） & 0.2392  & 0.2082  & 12.23\% & 0.0038  \\ 
        营业收入增长率（\%） & 0.2061  & 0.1584  & 9.30\% & 0.0044  \\ 
        净利润增长率（\%） & 0.2383  & 0.3639  & 21.36\% & -0.0268  \\ 
        总资产周转率（\%） & 0.2128  & 0.1993  & 11.70\% & 0.0016  \\ 
        应收账款周转率（\%） & 0.3371  & 0.1830  & 10.75\% & 0.0166 \\
      \bottomrule
    \end{tabular}
  \end{threeparttable}
\end{table}

在表\eqref{tab:indicator-hierarchy-different-data}中：调整额度的计算方式如下：

调整数值=（比亚迪信息效用-三家企业信息效用均值）*信息权重

通过表\eqref{tab:indicator-hierarchy-different-data}中的数据，参考三个上市的同行业公司还有2023年比亚迪财务报表，我们认为比亚迪2023年的最优资产负债率应该从65.73\%降低1.86\%，为63.87\%，略微低于行业资产负债率64.35\%。

综上所述，比亚迪2023年最优资产负债率区间为[63.87\%,65.73\%]。
\section{比亚迪最优化资本结构的建议}
\subsection{加强营运资金管理}
比亚迪公司在企业运营资本的管理方面需要进一步加强。需要根据宏观环境，新能源汽车行业发展阶段还有公司未来的战略规划来进行针对性的调整。这也意味着比亚迪需要指定符合企业未来战略目标的资本结构，不仅能适应未来多年的宏观环境，更加激烈的新能源行业竞争，还能满足日常企业运转的资金需要，满足企业长期投资发展的需要。  
%比亚迪公司在进一步增强自身的营运资金管理能力方面有着重要的举措。为提高对战略调整的敏感性和前瞻性，公司不仅需结合内外环境因素，也要考虑公司发展的阶段和新能源汽车行业的动态变化。这意味着比亚迪需要设计符合其发展阶段的融资结构，以确保既能支持日常运营需求，又能满足长期发展规划的资金需求。

另一方面，比亚迪还要注重金融风险的防范控制，努力将债务风险降到最低。为此，公司需要根据企业需求针对性地调整短期融资方案，根据企业发展情况和发展战略指定合适的融资计划，并在可控的风险下提高企业的价值，以优化企业的资本结构。这一综合性的资金管理策略不仅能够有效应对市场变化和风险挑战，还有助于确保公司资金稳健，并推动企业持续健康发展。
\subsection{合理分配债务期限结构}
其实通过之前分析比亚迪汽车2018年-2023年6年的财务数据和财务指标我们发现比亚迪汽车的资产负债率明显高于同行业同类型的企业，并且流动比率明显低于同行业公司。同时发现企业短期负债占比较高，长期负债占比较低的问题。如果不能通过期限匹配的方式进行融资，那么可能就会导致整体出现流动性风险。

详细来说，即企业应将债务的期限进行合理匹配。长期负债应用于长期项目，短期负债应用于短期项目。只有当资产期限比债务期限短时，资产能够产生足够的现金流来偿还债务，并促使企业在资产寿命终止时做好再投资的决策。这样不仅可以优化债务期限结构、降低企业不能及时偿债的财务风险，避免盲目投资。匹配的同时也要注意考虑企业能承担的风险和成本。
\subsection{合理调整企业融资渠道比重}
比亚迪新能源汽车公司在资本结构优化方面有一些待改进的地方，特别是在融资渠道的选择方面。目前，企业在融资方式方面主要采用债务融资，特别是大量使用短期债务进行融资，较少使用股权融资，使得融资来源比例失衡。为了达到最优的资本结构，比亚迪需要在未来调整股权融资在整个融资中的比重，更多地采用股权融资。
%为了实现最优资本结构，比亚迪亟需调整融资渠道，更多地利用股权融资渠道，如采用配股、增发等手段。

此外，从融资渠道来源来说，比亚迪的融资渠道越来越倾向于内源融资，外源融资占比越来越低。所以比亚迪在未来再次进行融资时，应该选取增发股票的方式。不仅可以显著降低资产负债率，改善资本结构，同时还可以改善资金来源占比失衡的问题。在选择增发股票这一举措时，可以考虑选择定向增发股票以保护现有股东的权益。 
%值得注意的是，比亚迪新能源汽车公司成立至今只进行过一次增发，而这次增发的效果明显改善了资本结构，资产负债率大幅下降，为公司的发展带来了积极影响。拉长比亚迪公司的成长线来看，保持原有股东权益不被改变，定向增发股票是一个能够有效缓解资金压力的思路，这点已在比亚迪公司目前的发展指标上得到了印证。

另一方面，比亚迪目前大量使用短期银行借款的融资方式虽然可以有效推动比亚迪短期业务发展，但是失衡的比例将会不利于比亚迪的长远发展和长期投资。因此，根据最优资本结构65\%建议，合理调整融资渠道将有助于比亚迪公司优化资本结构，降低财务风险，推动公司长期健康发展。
%另一方面，若融资方式只选择银行借款，虽然方便快捷，但会大大加重比亚迪公司偿还债务的压力，增大财务风险。因此，根据最优资本结构65\%建议，合理调整融资渠道将有助于比亚迪公司优化资本结构，降低财务风险，推动公司长期健康发展。
\chapter{结论}
\section{研究结论}
\section{研究展望}

%\chapter{测试}
% 其他部分
\backmatter
% 参考文献
\bibliography{reference/AccountThesis}  % 参考文献使用 BibTeX 编译(本科生)

% 研究生引入参考文献使用此版本
%\begin{kai}
%\bibliography{reference/refs}
%\end{kai}

% \printbibliography       % 参考文献使用 BibLaTeX 编译

% 致谢
%%
%% @Author           : Albert Wang
%% @Time             : 2023-01-08 21:21:02
%% @Description      :
%% @Email            : shadowofgost@outlook.com
%% @FilePath         : /Sudathesis-Bachelor-Template/contents/acknowledgements.tex
%% @LastTime         : 2023-01-10 11:07:29
%% @LastAuthor       : Albert Wang
%% @Software         : Vscode
%% @ Copyright Notice : Copyright (c) 2023 Albert Wang 王子睿, All Rights Reserved.
%%

% !TeX root = ../main.tex

\begin{acknowledgements}
当我写完这篇论文的时候已经是凌晨时分，我也意识到这是我本科学习生涯的结束。但同时，这也是我继续学习会计的起点。回忆过去的五年间本科学习的经历，内心充斥着各种各样的情绪。有开心、骄傲、自豪，也有在繁多的作业下的焦虑、煎熬和痛苦。在面对繁多的作业和课业压力的情况下，如何平衡好辅修课程、主修课程和各种竞赛便成为了一个极度困难的事情。曾经这对于我来说是一个走钢丝的行为，并且无数次打破了平衡变得焦头烂额。但是最终我都挺了过来直到站在了这里，站在了我无数次希望推开的门前，站在离我无数次希望拿到的学位前。我最终如尝所愿。

在完成整个毕业设计并进行撰写毕业论文的过程中，我逐渐体会到攥写论文本身并不是重点，如何按照正确的流程和方法拆分论文内容，查找组织数据并搭建模型才是重点，攥写论文不过是其中最为基础和关键的一步。

在这里我需要感谢指导老师龚蕾老师在我毕业设计中给与的指导，帮助我解决问题指明方向。感谢商学院的老师在我学习时提供的帮助，能帮助我从一个对会计金融管理等商科知识一无所知的小白逐渐成长为一个合格的财务人员 。同时感谢我的舍友在周末时能接受我的闹铃，帮助我及时起床上课学习会计的相关知识。

最后，我会在以后的日子中铭记这一段光辉灿烂又充满泪水的路程，在未来人生的漫漫征途中砥砺前行。

\end{acknowledgements}

%\chapter{附录}
% 附录
%\fancypagestyle{input}{%
%  \fancyhf{}%
%  \renewcommand\footrulewidth{0pt}%
%    \renewcommand\headrulewidth{0bp}%
%    \fancyfoot[C]{
%        \xiaowu
%      -\thepage-
%    }%
%}% 给下面引入的所有 pdf 标页码
%\includepdfset{pagecommand={\thispagestyle{input}}} %pdfpages宏包插入文件无页码的解决
%\setcounter{secnumdepth}{0}
%\includepdf[addtotoc={1,chapter,1,文献综述,cc,2,section,2,测试,aa},pages={1,2,3}]{./figures/statement.pdf}

% 本科生需要将附录放到声明之后，个人简历之前
%\appendix
%%%
%% @Author           : Albert Wang
%% @Time             : 2023-01-08 21:21:02
%% @Description      :
%% @Email            : shadowofgost@outlook.com
%% @FilePath         : /Sudathesis-Bachelor-Template/contents/appendix-survey.tex
%% @LastTime         : 2023-01-10 11:09:13
%% @LastAuthor       : Albert Wang
%% @Software         : Vscode
%% @ Copyright Notice : Copyright (c) 2023 Albert Wang 王子睿, All Rights Reserved.
%%

% !TeX root = ../main.tex

\begin{survey}
\label{cha:survey}

\title{Title of the Survey}
\maketitle


\tableofcontents


本科生的外文资料调研阅读报告。


\section{Figures and Tables}

\subsection{Figures}

An example figure in appendix (Figure~\ref{fig:appendix-survey-figure}).

\begin{figure}
  \centering
  \includegraphics[width=0.6\linewidth]{example-image-a.pdf}
  \caption{Example figure in appendix}
  \label{fig:appendix-survey-figure}
\end{figure}


\subsection{Tables}

An example table in appendix (Table~\ref{tab:appendix-survey-table}).

\begin{table}
  \centering
  \caption{Example table in appendix}
  \begin{tabular}{ll}
    \toprule
    File name       & Description                                         \\
    \midrule
    thuthesis.dtx   & The source file including documentaion and comments \\
    thuthesis.cls   & The template file                                   \\
    thuthesis-*.bst & BibTeX styles                                       \\
    thuthesis-*.bbx & BibLaTeX styles for bibliographies                  \\
    thuthesis-*.cbx & BibLaTeX styles for citations                       \\
    \bottomrule
  \end{tabular}
  \label{tab:appendix-survey-table}
\end{table}


\section{Equations}

An example equation in appendix (Equation~\eqref{eq:appendix-survey-equation}).
\begin{equation}
  \frac{1}{2 \uppi \symup{i}} \int_\gamma f = \sum_{k=1}^m n(\gamma; a_k) \mathscr{R}(f; a_k)
  \label{eq:appendix-survey-equation}
\end{equation}


\section{Citations}

Example citations in appendix.
\cite{abrahams99tex}
\cite{salomon1995advanced}
\cite{abrahams99tex,salomon1995advanced}


\bibliographystyle{unsrtnat}
\bibliography{ref/appendix}

\end{survey}
       % 本科生：外文资料的调研阅读报告
%%%
%% @Author           : Albert Wang
%% @Time             : 2023-01-08 21:21:02
%% @Description      :
%% @Email            : shadowofgost@outlook.com
%% @FilePath         : /Sudathesis-Bachelor-Template/contents/appendix-translation.tex
%% @LastTime         : 2023-01-10 11:09:24
%% @LastAuthor       : Albert Wang
%% @Software         : Vscode
%% @ Copyright Notice : Copyright (c) 2023 Albert Wang 王子睿, All Rights Reserved.
%%

% !TeX root = ../main.tex

\begin{translation}
\label{cha:translation}

\title{书面翻译题目}
\maketitle

\tableofcontents


本科生的外文资料书面翻译。


\section{图表示例}

\subsection{图}

附录中的图片示例（图~\ref{fig:appendix-translation-figure}）。

\begin{figure}
  \centering
  \includegraphics[width=0.6\linewidth]{example-image-a.pdf}
  \caption{附录中的图片示例}
  \label{fig:appendix-translation-figure}
\end{figure}


\subsection{表格}

附录中的表格示例（表~\ref{tab:appendix-translation-table}）。

\begin{table}
  \centering
  \caption{附录中的表格示例}
  \begin{tabular}{ll}
    \toprule
    文件名          & 描述                         \\
    \midrule
    thuthesis.dtx   & 模板的源文件，包括文档和注释 \\
    thuthesis.cls   & 模板文件                     \\
    thuthesis-*.bst & BibTeX 参考文献表样式文件    \\
    thuthesis-*.bbx & BibLaTeX 参考文献表样式文件  \\
    thuthesis-*.cbx & BibLaTeX 引用样式文件        \\
    \bottomrule
  \end{tabular}
  \label{tab:appendix-translation-table}
\end{table}


\section{数学公式}

附录中的数学公式示例（公式\eqref{eq:appendix-translation-equation}）。
\begin{equation}
  \frac{1}{2 \uppi \symup{i}} \int_\gamma f = \sum_{k=1}^m n(\gamma; a_k) \mathscr{R}(f; a_k)
  \label{eq:appendix-translation-equation}
\end{equation}


\section{文献引用}

文献引用示例\cite{abrahams99tex}。


\appendix

\section{附录}

附录的内容。


% 书面翻译的参考文献
\bibliographystyle{unsrtnat}
\bibliography{ref/appendix}

% 书面翻译对应的原文索引
\begin{translation-index}
  \nocite{salomon1995advanced}
  \bibliographystyle{unsrtnat}
  \bibliography{ref/appendix}
\end{translation-index}

\end{translation}
  % 本科生：外文资料的书面翻译

% 下面是引入附录
%%%
%% @Author           : Albert Wang
%% @Time             : 2023-04-06 22:03:52
%% @Description      :
%% @Email            : shadowofgost@outlook.com
%% @FilePath         : /sudathesis-soochow-university-latex-template/contents/appendix.tex
%% @LastTime         : 2023-12-25 19:06:19
%% @LastAuthor       : Albert Wang
%% @Software         : Vscode
%% @ Copyright Notice : Copyright (c) 2023 Albert Wang 王子睿, All Rights Reserved.
%%
% !TeX root = ../thuthesis-example.tex

\chapter{补充内容}

附录是与论文内容密切相关、但编入正文又影响整篇论文编排的条理和逻辑性的资料，例如某些重要的数据表格、计算程序、统计表等，是论文主体的补充内容，可根据需要设置。

附录中的图、表、数学表达式、参考文献等另行编序号，与正文分开，一律用阿拉伯数字编码，
但在数码前冠以附录的序号，例如“图~\ref{fig:appendix-figure}”，
“表~\ref{tab:appendix-table}”，“式\eqref{eq:appendix-equation}”等。


\section{插图}

% 附录中的插图示例（图~\ref{fig:appendix-figure}）。

\begin{figure}
  \centering
  \includegraphics[width=0.6\linewidth]{example-image-a.pdf}
  \caption{附录中的图片示例}
  \label{fig:appendix-figure}
\end{figure}


\section{表格}

% 附录中的表格示例（表~\ref{tab:appendix-table}）。

\begin{table}
  \centering
  \caption{附录中的表格示例}
  \begin{tabular}{ll}
    \toprule
    文件名          & 描述                         \\
    \midrule
    thuthesis.dtx   & 模板的源文件，包括文档和注释 \\
    thuthesis.cls   & 模板文件                     \\
    thuthesis-*.bst & BibTeX 参考文献表样式文件    \\
    thuthesis-*.bbx & BibLaTeX 参考文献表样式文件  \\
    thuthesis-*.cbx & BibLaTeX 引用样式文件        \\
    \bottomrule
  \end{tabular}
  \label{tab:appendix-table}
\end{table}


\section{数学表达式}

% 附录中的数学表达式示例（式\eqref{eq:appendix-equation}）。
\begin{equation}
  \frac{1}{2 \uppi \symup{i}} \int_\gamma f = \sum_{k=1}^m n(\gamma; a_k) \mathscr{R}(f; a_k)
  \label{eq:appendix-equation}
\end{equation}

% TODO！本科生请勿启用，会出现报错，硕博请自行评估是否需要来启用
\section{参考文献}

附录中的参考文献示例（\cite{carlson1981two} 和 \cite{carlson1981two,taylor1983scanning,taylor1981study}）。

%\printbibliography


% 致谢
%%%
%% @Author           : Albert Wang
%% @Time             : 2023-01-08 21:21:02
%% @Description      :
%% @Email            : shadowofgost@outlook.com
%% @FilePath         : /Sudathesis-Bachelor-Template/contents/acknowledgements.tex
%% @LastTime         : 2023-01-10 11:07:29
%% @LastAuthor       : Albert Wang
%% @Software         : Vscode
%% @ Copyright Notice : Copyright (c) 2023 Albert Wang 王子睿, All Rights Reserved.
%%

% !TeX root = ../main.tex

\begin{acknowledgements}
当我写完这篇论文的时候已经是凌晨时分，我也意识到这是我本科学习生涯的结束。但同时，这也是我继续学习会计的起点。回忆过去的五年间本科学习的经历，内心充斥着各种各样的情绪。有开心、骄傲、自豪，也有在繁多的作业下的焦虑、煎熬和痛苦。在面对繁多的作业和课业压力的情况下，如何平衡好辅修课程、主修课程和各种竞赛便成为了一个极度困难的事情。曾经这对于我来说是一个走钢丝的行为，并且无数次打破了平衡变得焦头烂额。但是最终我都挺了过来直到站在了这里，站在了我无数次希望推开的门前，站在离我无数次希望拿到的学位前。我最终如尝所愿。

在完成整个毕业设计并进行撰写毕业论文的过程中，我逐渐体会到攥写论文本身并不是重点，如何按照正确的流程和方法拆分论文内容，查找组织数据并搭建模型才是重点，攥写论文不过是其中最为基础和关键的一步。

在这里我需要感谢指导老师龚蕾老师在我毕业设计中给与的指导，帮助我解决问题指明方向。感谢商学院的老师在我学习时提供的帮助，能帮助我从一个对会计金融管理等商科知识一无所知的小白逐渐成长为一个合格的财务人员 。同时感谢我的舍友在周末时能接受我的闹铃，帮助我及时起床上课学习会计的相关知识。

最后，我会在以后的日子中铭记这一段光辉灿烂又充满泪水的路程，在未来人生的漫漫征途中砥砺前行。

\end{acknowledgements}


% 声明
%\statement
% 将签字扫描后的声明文件 scan-statement.pdf 替换原始页面
% \statement[file=scan-statement.pdf]
% 本科生编译生成的声明页默认不加页脚，插入扫描版时再补上；
% 研究生编译生成时有页眉页脚，插入扫描版时不再重复。
% 也可以手动控制是否加页眉页脚
% \statement[page-style=empty]
% \statement[file=scan-statement.pdf, page-style=plain]

% 个人简历、在学期间完成的相关学术成果
% 本科生可以附个人简历，也可以不附个人简历
%%%
%% @Author           : Albert Wang
%% @Time             : 2023-01-08 21:21:02
%% @Description      :
%% @Email            : shadowofgost@outlook.com
%% @FilePath         : /Sudathesis-Bachelor-Template/contents/resume.tex
%% @LastTime         : 2023-01-10 11:11:42
%% @LastAuthor       : Albert Wang
%% @Software         : Vscode
%% @ Copyright Notice : Copyright (c) 2023 Albert Wang 王子睿, All Rights Reserved.
%%

% !TeX root = ../main.tex

\begin{resume}

  \section*{个人简历}

  197× 年 ×× 月 ×× 日出生于四川××县。

  1992 年 9 月考入××大学化学系××化学专业，1996 年 7 月本科毕业并获得理学学士学位。

  1996 年 9 月免试进入清华大学化学系攻读××化学博士至今。


  \section*{在学期间完成的相关学术成果}

  \subsection*{学术论文}

  \begin{achievements}
    \item Yang Y, Ren T L, Zhang L T, et al. Miniature microphone with silicon-based ferroelectric thin films[J]. Integrated Ferroelectrics, 2003, 52:229-235.
    \item 杨轶, 张宁欣, 任天令, 等. 硅基铁电微声学器件中薄膜残余应力的研究[J]. 中国机械工程, 2005, 16(14):1289-1291.
    \item 杨轶, 张宁欣, 任天令, 等. 集成铁电器件中的关键工艺研究[J]. 仪器仪表学报, 2003, 24(S4):192-193.
    \item Yang Y, Ren T L, Zhu Y P, et al. PMUTs for handwriting recognition. In press[J]. (已被Integrated Ferroelectrics录用)
  \end{achievements}


  \subsection*{专利}

  \begin{achievements}
    \item 任天令, 杨轶, 朱一平, 等. 硅基铁电微声学传感器畴极化区域控制和电极连接的方法: 中国, CN1602118A[P]. 2005-03-30.
    \item Ren T L, Yang Y, Zhu Y P, et al. Piezoelectric micro acoustic sensor based on ferroelectric materials: USA, No.11/215, 102[P]. (美国发明专利申请号.)
  \end{achievements}

\end{resume}


% 指导教师/指导小组评语
% 本科生不需要
%%%
%% @Author           : Albert Wang
%% @Time             : 2023-01-08 21:21:02
%% @Description      :
%% @Email            : shadowofgost@outlook.com
%% @FilePath         : /Sudathesis-Bachelor-Template/contents/comments.tex
%% @LastTime         : 2023-01-10 11:11:18
%% @LastAuthor       : Albert Wang
%% @Software         : Vscode
%% @ Copyright Notice : Copyright (c) 2023 Albert Wang 王子睿, All Rights Reserved.
%%

% !TeX root = ../main.tex

\begin{comments}
% \begin{comments}[name = {指导小组评语}]
% \begin{comments}[name = {Comments from Thesis Supervisor}]
% \begin{comments}[name = {Comments from Thesis Supervision Committee}]

  论文提出了……

\end{comments}


% 答辩委员会决议书
% 本科生不需要
%%%
%% @Author           : Albert Wang
%% @Time             : 2023-01-08 21:21:02
%% @Description      :
%% @Email            : shadowofgost@outlook.com
%% @FilePath         : /Sudathesis-Bachelor-Template/contents/resolution.tex
%% @LastTime         : 2023-01-10 11:11:37
%% @LastAuthor       : Albert Wang
%% @Software         : Vscode
%% @ Copyright Notice : Copyright (c) 2023 Albert Wang 王子睿, All Rights Reserved.
%%

% !TeX root = ../main.tex

\begin{resolution}

  论文提出了……

  论文取得的主要创新性成果包括：

  1. ……

  2. ……

  3. ……

  论文工作表明作者在×××××具有×××××知识，具有××××能力，论文××××，答辩××××。

  答辩委员会表决，（×票/一致）同意通过论文答辩，并建议授予×××（姓名）×××（门类）学博士/硕士学位。

\end{resolution}


% 本科生的综合论文训练记录表（扫描版）
% \record{file=scan-record.pdf}

\end{document}
